\chapter{Pediatrics}

\medterm{Bilirubin Encephalopathy}-- Neurologic injury caused by deposition of unconjugated bilirubin in the brain, especially the basal ganglia. Severe neonatal hyperbilirubinemia can produce acute lethargy, abnormal tone, feeding difficulty, seizures, and permanent kernicterus.

\medterm{Ganglioneuroblastoma}-- A pediatric tumor of the sympathetic nervous system containing both neuroblastic and ganglion-cell elements. It most often arises in the adrenal or paraspinal region and may present with an abdominal mass, pain, hypertension, or systemic symptoms.

\medterm{Natal Teeth}-- Teeth present at birth, most often involving the lower central incisors. They may be loose, interfere with feeding, or injure the infant's tongue and require dental assessment.

\medterm{Neonatal Conjunctivitis}-- Inflammation of the conjunctiva during the first month of life, caused by chemical irritation, bacteria, or viruses acquired during delivery. Purulent discharge, eyelid swelling, and corneal involvement require urgent evaluation.

\medterm{Infant Botulism}-- A neuromuscular illness caused by intestinal colonization with Clostridium botulinum and production of botulinum neurotoxin in infants. It causes constipation, poor feeding, weak cry, hypotonia, facial weakness, and descending paralysis.

\medterm{Encopresis}-- Recurrent passage of stool in inappropriate places by a child who is developmentally old enough for bowel control. It is commonly associated with chronic constipation and fecal retention, with overflow soiling around impacted stool.

\medterm{Intestinal Malrotation}-- Congenital failure of normal midgut rotation and fixation, predisposing to a narrow mesenteric base, volvulus, and duodenal obstruction. Infants may present with bilious vomiting, abdominal distension, or bowel ischemia; sudden symptoms require emergency assessment. Upper gastrointestinal contrast imaging or ultrasound supports diagnosis, and symptomatic malrotation is treated surgically.

\medterm{Acne}-- A chronic inflammatory disorder of the pilosebaceous unit, extremely common during adolescence. Pathogenesis involves sebum production, follicular hyperkeratinization, Cutibacterium acnes colonization, and inflammation. Grading: comedonal, papulopustular, nodulocystic. Treatment: topical retinoids, benzoyl peroxide, and antibiotics for mild to moderate; isotretinoin for severe or refractory disease. Early treatment prevents scarring.

\medterm{Acute Otitis Media}-- Infection of the middle ear space with acute onset of symptoms and signs of middle ear inflammation. Diagnosis requires middle ear effusion plus acute symptoms (ear pain, fever, irritability) or signs (bulging tympanic membrane, erythema). Most common pathogen: Streptococcus pneumoniae. Treatment: amoxicillin for 10 days; amoxicillin-clavulanate for treatment failure or concurrent sinusitis.

\medterm{ADHD}-- Attention deficit hyperactivity disorder: a neurodevelopmental disorder with inattention, hyperactivity, and impulsivity inappropriate for age, affecting function in school/home; onset before 12 years, persistence >=6 months in >=2 settings. Subtypes: inattentive, hyperactive-impulsive, combined. Comorbidities: learning disorders, anxiety, depression, ODD, tics. Diagnosis: clinical (DSM-5 criteria), rating scales (Conners, Vanderbilt), history from parents/teachers; exclude hearing/vision, sleep, anxiety, and trauma. Management: behavioral/parent training first (preschool), stimulants (methylphenidate) or non-stimulants (atomoxetine, guanfacine) for school-age, structured environment, and educational support; monitor growth and side effects.

\medterm{Adolescent Health}-- The health and wellbeing of young people (10-19 years), addressing the physical, psychological, and social transitions of puberty and adolescence. Key issues: growth/development (puberty, menarche), nutrition (obesity, eating disorders), mental health (depression, anxiety, suicide—leading causes of morbidity), sexual health (contraception, STIs, pregnancy), substance use (tobacco, alcohol, drugs), injuries/violence, and preventive care (immunizations—HPV, Tdap; screening). Confidentially provided care, HEADSSS assessment (Home, Education, Activities, Drugs, Sexuality, Suicide, Safety), and a non-judgmental approach are essential. Early identification and intervention improve long-term outcomes.

\medterm{Alagille Syndrome}-- An autosomal dominant disorder caused by mutations in JAG1 (95\%) or NOTCH2 (<5\%), affecting the Notch signaling pathway. Features: chronic cholestasis from paucity of intrahepatic bile ducts (90-100\%), cardiac anomalies (peripheral pulmonary stenosis most common, also tetralogy of Fallot), butterfly vertebrae (50\%), posterior embryotoxon (a ring in the anterior chamber of the eye, 90\%), and characteristic facies (broad forehead, deep-set eyes, pointed chin, 90-95\%). Presents in infancy with prolonged jaundice, pruritus, and failure to thrive. Other findings: renal anomalies (30-50\%), hypercholesterolemia, and coagulopathy from fat malabsorption. Diagnosis: clinical criteria (Alagille score) and JAG1/NOTCH2 genetic testing. Management: ursodeoxycholic acid for cholestasis, fat-soluble vitamin supplementation, nutritional support, cardiac repair as needed, and liver transplantation for end-stage liver disease (15-20\%).

\medterm{Angelman Syndrome}-- A genetic disorder caused by loss of function of the UBE3A gene on chromosome 15 (maternal deletion or paternal uniparental disomy). Features: severe intellectual disability, absent speech, seizures, ataxia, and frequent laughter. Characteristic happy demeanor with hand-flapping movements. Management: seizure control, physical therapy, and communication support.

\medterm{Annular Pancreas}-- A congenital anomaly where a ring of pancreatic tissue completely or partially encircles the second portion of the duodenum, causing duodenal obstruction. It results from failure of the ventral pancreatic bud to completely rotate around the duodenum during embryogenesis. Associated with Down syndrome, duodenal atresia (coexisting in 50\%), and other GI anomalies. Presents in the newborn period with bilious or non-bilious vomiting (depending on the position relative to the ampulla) and polyhydramnios (prenatal). In adults, may present with recurrent pancreatitis or peptic ulcer disease. Diagnosis: upper GI series shows a narrowed, eccentric duodenum with obstructive findings; CT or MRCP demonstrates the pancreatic tissue ring encasing the duodenum. Treatment: surgical bypass (duodeno-duodenostomy or duodeno-jejunostomy) without resection of the annular tissue (risk of pancreatitis and fistula). Prognosis: excellent after surgical repair.

\medterm{Anorectal Malformation}-- A spectrum of congenital anomalies involving the anus and rectum, classified by the location of the rectal pouch relative to the puborectalis sling: Low (passes through the levator muscle, good prognosis) — perineal fistula, vestibular fistula, or covered anus. High (ends above the levator muscle, poor prognosis, 60\%) — rectourethral fistula (boys), rectovaginal fistula (girls), or rectobladder neck fistula. Cloaca (girls): a single common channel for the urethra, vagina, and rectum. Presents with absence of a normal anal opening on perineal examination and failure to pass meconium within 24 hours. Associated anomalies (VACTERL workup): vertebral, cardiac (VSD), tracheoesophageal, renal (VUR, renal agenesis), and limb. Diagnosis: invertogram (Wangensteen-Rice X-ray) and cross-table lateral X-ray to determine the level. MRI for complex cases. Treatment: low lesions — primary perineal anoplasty. High lesions — colostomy in the newborn period, followed by posterior sagittal anorectoplasty (PSARP, deVries-Peña procedure) at 3-6 months. Prognosis: low lesions — excellent bowel control. High lesions — variable fecal continence, requiring bowel management programs (enemas, laxatives, antegrade continence enema).

\medterm{Anorexia Nervosa}-- An eating disorder characterized by restriction of energy intake, intense fear of weight gain, and distorted body image. Highest mortality of any psychiatric disorder. Features: emaciation, lanugo, bradycardia, hypotension, amenorrhea, and electrolyte abnormalities. Diagnosis by DSM-5 criteria. Treatment: nutritional rehabilitation, cognitive behavioral therapy, and family-based therapy. Medical stabilization when severely malnourished.

\medterm{Apert Syndrome}-- An autosomal dominant craniosynostosis syndrome caused by FGFR2 mutations (most commonly Ser252Trp or Pro253Arg). Features: craniosynostosis (bilateral coronal, resulting in turribrachycephaly — tower-shaped skull), midface hypoplasia, proptosis, hypertelorism, and syndactyly of the hands and feet (mitten or spade-like hands, complete fusion of digits). Cleft palate occurs in 30\%. May have intellectual disability (variable). Other associations: cervical spine fusion, congenital heart disease, and hydrocephalus. Diagnosis: clinical and genetic testing. Management: staged surgical reconstruction — cranial vault remodeling (3-6 months), midface advancement (Le Fort III, 6-12 years), and syndactyly release (6-18 months). Lifelong neurosurgical and dental monitoring.

\medterm{Apgar Score}-- Rapid clinical assessment of newborn well-being performed at 1 and 5 minutes after birth. Five criteria scored 0-2 each: Appearance (skin color), Pulse (heart rate), Grimace (reflex irritability), Activity (muscle tone), and Respiration. Total score 0-10. Score 7-10 normal; 4-6 moderately depressed; 0-3 severely depressed. Determines need for immediate resuscitation. May extend to 10 and 20 minutes in severely depressed infants.

\medterm{Apnea of Prematurity}-- Cessation of breathing for more than 20 seconds or shorter episodes with bradycardia or desaturation in infants less than 34-37 weeks gestation. Three types: central (no respiratory effort), obstructive (airway collapse), and mixed. Treatment: caffeine citrate (first-line), continuous positive airway pressure, and methylxanthine therapy. Monitor with cardiorespiratory monitors. Most resolve by 34-44 weeks corrected gestational age.

\medterm{Argininosuccinic Aciduria}-- Autosomal recessive urea cycle disorder caused by deficiency of argininosuccinate lyase (ASL gene, 7cen-q11.2). Biochemical hallmark: accumulation of argininosuccinic acid in blood, urine, and CSF with hyperammonemia, low citrulline, and elevated argininosuccinic acid. Neonatal form presents within days of birth with poor feeding, vomiting, lethargy, tachypnea, seizures, and progressing to coma; milder late-onset forms present later with developmental delay, intellectual disability, hepatic dysfunction, and brittle hair (trichorrhexis nodosa). Diagnosis: plasma amino acids (elevated citrulline, argininosuccinic acid), urine orotic acid, enzyme assay in erythrocytes, and ASL genetic testing. Treatment: ammonia scavenging (sodium benzoate, sodium phenylbutyrate, glycerol phenylbutyrate), arginine supplementation (provides an alternative excretion route via argininosuccinate in urine), low-protein diet, and liver transplantation in severe cases.

\medterm{Arnold-Chiari Malformation}-- Herniation of cerebellar tonsils through the foramen magnum into the upper cervical spinal canal. Type II (myelomeningocele-associated) is the most common form in pediatrics. Associated with hydrocephalus, syringomyelia, and spinal cord dysfunction. Symptoms include headache, neck pain, and lower cranial nerve palsies. Treatment: posterior fossa decompression with duraplasty. Often coexists with myelomeningocele requiring concurrent repair.

\medterm{Babinski Reflex}-- An extensor plantar response present from birth to approximately 2 years. When the lateral sole is stroked, the great toe dorsiflexes and other toes fan outward. Disappears as the corticospinal tract myelinates. Persistence beyond age 2 may indicate neurologic abnormality. Pathologic in older children and adults, suggesting upper motor neuron lesion. One of the primitive reflexes assessed in infants.

\medterm{Bacterial Meningitis}-- Infection of the meninges by bacteria, a medical emergency. In neonates: Group B Streptococcus, E. coli, Listeria. In older children: Neisseria meningitidis, Streptococcus pneumoniae. Presents with fever, headache, neck stiffness, altered consciousness, and Kernig/Brudzinski signs. Diagnosis by lumbar puncture. Treatment: immediate empiric antibiotics (ceftriaxone and vancomycin) without delaying for imaging.

\medterm{Ballard Score}-- Clinical assessment of gestational age in newborns using 6 physical and 6 neuromuscular maturity criteria. Physical parameters: skin texture, ear form and recoil, eye appearance, breast tissue, and genitalia. Neuromuscular: posture, arm and leg recoil, popliteal angle, scarf sign, and heel-to-ear. Each scored 0-5; total converted to gestational age in weeks. Most accurate near term; used when prenatal dating is uncertain.

\medterm{Beckwith-Wiedemann Syndrome}-- An overgrowth disorder associated with 11p15.5 imprinting defects. Features: macrosomia, macroglossia, omphalocele, visceromegaly, and neonatal hypoglycemia. Increased risk of embryonal tumors (Wilms tumor, hepatoblastoma). Management: glucose monitoring, tumor surveillance with abdominal ultrasound and AFP levels, and surgical correction of structural anomalies.

\medterm{Bilious Vomiting}-- Vomiting of green/yellow, bilious material in the newborn or infant, defined as emerald green or yellow vomitus that contains bile. In neonates (first days to weeks), it is a surgical emergency until proven otherwise, often indicating intestinal obstruction at or below the ampulla of Vater: differential includes duodenal atresia (Down syndrome, polyhydramnios, "double bubble" radiograph, bilious vomiting on day 1), malrotation with midgut volvulus (true emergency, "corkscrew" upper GI series, Ladd procedure), jejunoileal atresia (distension, bilious vomitus, intraperitoneal calcifications on radiograph), meconium ileus (cystic fibrosis), Hirschsprung disease, and necrotizing enterocolitis. In older children: gastroenteritis, intussusception (currant-jelly stool), appendicitis, post-treatment vomiting due to pyloric stenosis (typically NON-bilious). Workup: abdominal X-ray, contrast upper GI series, ultrasound, electrolyte/fluid resuscitation, NPO, NG decompression, surgical consultation.

\medterm{Binge Eating Disorder}-- Recurrent episodes of eating large quantities of food with a sense of lack of control, without regular compensatory behavior. Associated with obesity, depression, and anxiety. Diagnosis by DSM-5 criteria. Treatment: cognitive behavioral therapy, interpersonal therapy, and SSRIs. Weight management with dietary counseling and exercise.

\medterm{Biotinidase Deficiency}-- An autosomal recessive inborn error of biotin recycling caused by BTD gene mutations, resulting in inability to cleave biotin from biocytin (proteolytic degradation product of biotin-dependent carboxylases). This leads to multiple carboxylase deficiency affecting four biotin-dependent enzymes: pyruvate carboxylase, propionyl-CoA carboxylase, 3-methylcrotonyl-CoA carboxylase, and acetyl-CoA carboxylase. Presents between 1 week and 10 years of age (mean 3 months). Features: neurological — seizures (myoclonic, infantile spasms), hypotonia, ataxia, hearing loss, and optic atrophy; cutaneous — periorificial dermatitis, alopecia, and candidal infections; metabolic — metabolic acidosis, ketolactic acidosis, and organic aciduria (3-hydroxyisovalerate, 3-methylcrotonylglycine, lactate, propionate). Diagnosis: biotinidase enzyme assay in serum and BTD genetic testing. Included in newborn screening in many countries. Treatment: oral biotin (5-20 mg/day) — a dramatically effective and inexpensive therapy. Starting biotin promptly reverses most symptoms except sensorineural hearing loss and optic atrophy, which may be irreversible if delayed.

\medterm{Blue Baby}-- Generic term for an infant with cyanosis (bluish skin, lips, nail beds) from arterial oxygen desaturation. Differential: cyanotic congenital heart disease (right-to-left shunting; classic "5 T's" of Fallot, transposition of the great arteries, tricuspid atresia, truncus arteriosus, total anomalous pulmonary venous connection; plus hypoplastic left heart syndrome, pulmonary atresia), persistent pulmonary hypertension of the newborn, respiratory disorders (respiratory distress syndrome, meconium aspiration, congenital diaphragmatic hernia), methemoglobinemia (reduced hemoglobin cannot bind O2), polycythemia, and sepsis. Differentiate central cyanosis (warm, blue tongue/lips from low arterial O2) from peripheral (cold, acrocyanosis, normal O2 saturation). Workup: hyperoxia test (100\% FiO2 for 10 min; $\mathrm{PaO_2} > 150$ mmHg suggests pulmonary cause; $< 100$ mmHg suggests cyanotic CHD), pre/post-ductal saturations, chest X-ray, echocardiography, ABG/anion gap, sepsis workup. Emergency management: prostaglandin E1 infusion to maintain ductus arteriosus in duct-dependent lesions, transfer to cardiac center.

\medterm{Bochdalek Hernia}-- A congenital diaphragmatic hernia (CDH) occurring through the posterolateral foramen of Bochdalek, resulting from failure of the pleuroperitoneal membrane to close by the 8th week of gestation. The left side is affected in 85\% of cases (the right pleuroperitoneal canal closes earlier). Abdominal organs (stomach, spleen, small intestine, colon, and sometimes the left lobe of the liver) herniate into the chest, causing pulmonary hypoplasia and pulmonary hypertension. Presents at birth with severe respiratory distress, cyanosis, scaphoid abdomen, and decreased breath sounds on the affected side with bowel sounds heard in the chest. Diagnosis: prenatal ultrasound (polyhydramnios, mediastinal shift, herniated organs), postnatal chest X-ray shows abdominal organs in the thorax. Treatment: endotracheal intubation (avoid bag-mask ventilation — distends the stomach), nasogastric decompression, and surgical repair via an abdominal approach once pulmonary hypertension is managed. Delayed surgical repair after physiologic stabilization improves outcomes. High-frequency ventilation, inhaled nitric oxide, and ECMO for severe pulmonary hypertension. Prognosis depends on the degree of pulmonary hypoplasia; LHR (lung-to-head ratio) is the best prenatal predictor.

\medterm{Brachial Plexus Injury}-- Damage to the brachial plexus nerves during birth, most commonly from shoulder dystocia or excessive traction. Erb's palsy (C5-C6) presents with waiter's tip position (arm adducted, elbow extended, forearm pronated). Klumpke's palsy (C8-T1) affects the hand and may cause Horner syndrome. Most cases resolve spontaneously. Physical therapy initiated early; surgical exploration considered if no recovery by 3-6 months.

\medterm{Bradycardia of Prematurity}-- Heart rate less than 100 beats per minute in a premature infant, often associated with apnea and oxygen desaturation. May be caused by vagal stimulation, reflux, or immaturity of the cardiac conduction system. Treatment: gentle stimulation, caffeine, and addressing underlying causes. Usually self-limited and resolves as the infant matures. Monitoring with continuous cardiorespiratory monitoring.

\medterm{Branchial Cleft Cyst}-- Congenital cystic mass arising from incomplete obliteration of branchial cleft apparatus during embryogenesis. Most are second branchial cleft (95\%), presenting in childhood or adolescence as a painless, fluctuant anterior neck mass along the anterior border of the sternocleidomastoid (typically at the junction of upper and middle thirds), inferior to the angle of the mandible; may enlarge with upper respiratory infections; fistula to skin can discharge; tract passes between internal and external carotid arteries to the tonsillar fossa. First branchial cleft cysts (type I or II) occur near the ear or angle of mandible, and may present with otorrhea. Third/fourth branchial anomalies present as left thyroid or pyriform sinus fistulas. Diagnosis: ultrasound, contrast fistulography, CT, MRI. Treatment: complete surgical excision of cyst and tract (after acute infection resolves); incomplete excision risks recurrence.

\medterm{Breast Engorgement}-- See obstetrics and gynecology.

\medterm{Breath-Holding Spells}-- A benign paroxysmal phenomenon in young children (6 months to 6 years), provoked by frustration, pain, or fear. Cyanotic type (more common): child cries, forcefully exhales, becomes cyanotic, and may lose consciousness briefly. Pallid type: triggered by sudden pain/fright, child becomes pale and limp, may have brief seizure-like activity due to reflex anoxic seizure. Pathophysiology: autonomic dysregulation, vagal overactivity. Diagnosis: clinical; EEG normal (distinguishing from epilepsy). Treatment: parental reassurance, management of triggers, iron supplementation if anemic. Spells resolve spontaneously by age 6. No increased risk of epilepsy or other neurological problems.

\medterm{Bronchiolitis}-- Lower respiratory tract infection causing inflammation and obstruction of small airways. Most common cause is respiratory syncytial virus. Affects infants under 12 months. Presents with wheezing, tachypnea, nasal flaring, and retractions. Diagnosis is clinical. Treatment: supportive care with nasal suctioning, supplemental oxygen, and respiratory support. Severe cases may require CPAP or mechanical ventilation. No effective antiviral therapy.

\medterm{Bronchopulmonary Dysplasia}-- A chronic lung disease of premature infants, resulting from lung injury from mechanical ventilation, oxygen toxicity, and inflammation. Most commonly occurs in infants born at less than 28 weeks gestation with birth weight less than 1500 grams. Pathophysiology involves arrested alveolar and vascular development with simplified alveolar structures, decreased alveolar number, and pulmonary vascular remodeling. The National Institutes of Health consensus definition classifies severity at 36 weeks postmenstrual age or at discharge: mild (no supplemental oxygen), moderate (oxygen less than 30\%), or severe (oxygen at least 30\% or positive pressure ventilation). Pathogenesis is multifactorial: ventilator-induced lung injury, oxygen toxicity, infection, patent ductus arteriosus, and fluid overload. Management: lung-protective ventilation strategies, judicious oxygen therapy targeting saturations 90-95\%, fluid restriction, diuretics, caffeine, and systemic corticosteroids for severe cases not weaning from ventilation. Nutritional support with high-calorie formulas to promote lung growth. Long-term complications include pulmonary hypertension, reactive airway disease, and neurodevelopmental impairment. Most infants improve with lung growth, but pulmonary function abnormalities may persist into adulthood.

\medterm{Bulimia Nervosa}-- An eating disorder characterized by recurrent episodes of binge eating followed by compensatory behaviors (purging, laxatives, excessive exercise). Self-esteem overly influenced by body shape and weight. Complications: electrolyte abnormalities (hypokalemia, metabolic alkalosis), dental erosion, esophageal tears, and parotid gland enlargement. Treatment: cognitive behavioral therapy and SSRIs (fluoxetine).

\medterm{Caput Succedaneum}-- Edema of the scalp presenting as a soft, boggy swelling that crosses suture lines. Caused by pressure of the vaginal walls on the fetal head during labor. Resolves within hours to days. Differentiated from cephalohematoma (subperiosteal, does not cross suture lines) and subgaleal hemorrhage (life-threatening, crosses suture lines). No treatment required; spontaneous resolution.

\medterm{Cerebral Palsy}-- A group of permanent disorders of movement and posture development causing activity limitation, attributed to non-progressive disturbances occurring in the developing fetal or infant brain. It is the most common motor disability in childhood, affecting approximately 1 in 500 live births. Classification by motor type: spastic (80 percent, most common, includes hemiplegia, diplegia, and quadriplegia), dyskinetic (athetoid, choreoathetoid, dystonic), ataxic, hypotonic, and mixed. Spastic cerebral palsy is subclassified by distribution: hemiplegia (unilateral arm and leg), diplegia (lower limbs more affected, common in premature infants with periventricular leukomalacia), and quadriplegia (all four limbs, often with intellectual disability and seizures). Causes: prematurity, intrauterine infection, perinatal hypoxia-ischemia, intracranial hemorrhage, and genetic abnormalities. Diagnosis is clinical, based on delayed motor milestones, abnormal muscle tone, and persistent primitive reflexes. Management is multidisciplinary: physical therapy, occupational therapy, speech therapy, orthotics, spasticity management (botulinum toxin, oral baclofen, intrathecal baclofen pump), orthopaedic surgery for contractures, and management of associated conditions including epilepsy, feeding difficulties, and learning disabilities. Prognosis depends on severity and associated impairments; most children survive to adulthood with appropriate support.

\medterm{CHARGE Syndrome}-- A genetic disorder caused by CHD7 gene mutations (autosomal dominant, most cases de novo) affecting multiple organ systems. The mnemonic CHARGE describes: Coloboma (coloboma of the eye, 80-90\%), Heart defects (tetralogy of Fallot, VSD, ASD, PDA, 75-85\%), Atresia choanae (choanal atresia/stenosis, 50-60\%), Retarded growth and development (growth hormone deficiency, intellectual disability), Genital hypoplasia (cryptorchidism, micropenis, 50-60\% in males), and Ear anomalies (characteristic cup-shaped ears with prominent antihelix, hearing loss, 90-95\%). Presents with life-threatening respiratory distress at birth from choanal atresia. Other common features: facial palsy, cleft lip/palate, and feeding difficulties. Diagnosis: clinical criteria (Verloes or Blake) and CHD7 genetic testing. Management: multidisciplinary — airway management, cardiac repair, hearing aids/cochlear implants, growth hormone therapy, and developmental support.

\medterm{Chickenpox}-- Primary infection with Varicella-Zoster Virus, characterized by pruritic, vesicular rash in different stages (macules, papules, vesicles, crusts). Highly contagious via respiratory droplets and contact with lesions. Incubation 10-21 days. Complications include secondary bacterial infection, pneumonia, encephalitis, and Reye syndrome. Prevention by varicella vaccine. Treatment: supportive care; acyclovir for severe cases.

\medterm{Child Abuse}-- Physical, sexual, or emotional harm or neglect of a child by a caregiver. Recognized by inconsistent histories, injury patterns inconsistent with mechanism, and behavioral indicators. Types: physical abuse, sexual abuse, neglect, and emotional abuse. Mandatory reporting requirements for healthcare providers. Evaluation includes full skeletal survey for suspected physical abuse, laboratory tests, and social work involvement.

\medterm{Child Neglect}-- The failure of a caregiver to provide adequate food, shelter, clothing, supervision, education, or medical care. The most common form of child maltreatment. Types: physical neglect, educational neglect, emotional neglect, and medical neglect. May result in failure to thrive, developmental delay, and chronic illness. Identified through failure to attend appointments, untreated medical conditions, and inadequate living conditions.

\medterm{Chlamydia}-- The most common reportable bacterial STI, caused by Chlamydia trachomatis. Often asymptomatic in women, causing cervicitis, urethritis, and pelvic inflammatory disease. Complications include tubal factor infertility and ectopic pregnancy. Diagnosis by nucleic acid amplification test. Treatment: azithromycin or doxycycline. Partner treatment essential to prevent reinfection. Screening recommended for all sexually active women under 25.

\medterm{Choanal Atresia}-- A congenital obstruction of the posterior nasal passage (choana) due to failure of the bucconasal membrane to rupture during embryogenesis. Bilateral choanal atresia presents as a respiratory emergency in the newborn, who is an obligate nasal breather — cyclic cyanosis that improves with crying. Unilateral atresia presents later with persistent unilateral nasal discharge and obstruction. 70\% are bony, 30\% are membranous. Associated with CHARGE syndrome (coloboma, heart defects, choanal atresia, retardation, genital hypoplasia, ear anomalies) and other anomalies. Diagnosis: inability to pass a nasogastric tube through the nose, confirmed by CT. Treatment: immediate airway management (oral airway, McGovern nipple), followed by surgical repair via transnasal endoscopic choanoplasty with stenting. Unilateral cases are repaired electively.

\medterm{Clavicle Fracture}-- The most common birth-related fracture, occurring in approximately 1-2\% of vaginal deliveries. Risk factors include shoulder dystocia, macrosomia, and breech presentation. Presents with asymmetric arm movements, palpable crepitus, and reluctance to move the affected arm. Diagnosis by X-ray. Treatment: arm immobilization against the body (pinned diaper or arm sling). Heals rapidly in newborns with excellent prognosis.

\medterm{Cleft Lip}-- A congenital malformation: incomplete fusion of the lip (and often palate) during embryologic development, causing a facial cleft; cleft lip and/or cleft palate are the most common craniofacial birth defects. Features: cleft lip (unilateral/bilateral) and/or cleft palate (hard/soft); feeding difficulties, speech problems, recurrent otitis media/hearing loss, dental anomalies, and psychosocial impact; associated with syndromes (Pierre Robin, Van der Woude, trisomy 13/18). Diagnosis: prenatal ultrasound; clinical at birth. Management: multidisciplinary team—feeding support (special nipples), surgical repair (lip ~3-6 months, palate ~9-18 months), speech therapy, audiology/ENT, dental/orthodontics, and psychologic support.

\medterm{Cloacal Exstrophy}-- The most severe congenital anomaly of the ventral abdominal wall, also called the OEIS complex (Omphalocele, Exstrophy of the bladder, Imperforate anus, Spinal defects). Incidence 1 in 200,000-400,000 births. Features: an exstrophic central bowel segment (cecum, between two hemi-bladders), omphalocele, imperforate anus, and spinal anomalies (myelomeningocele, tethered cord). Girls have a duplicated vagina and uterus; boys have an epispadiac, short phallus. There is near-universal association with Chiari malformation and hydrocephalus. Diagnosis: prenatal ultrasound shows a large midline abdominal wall defect with cystic structures (hemi-bladders) and spinal abnormalities. Treatment: staged surgical reconstruction — immediate bowel coverage (silo), separation of bowel from bladder, closure of the abdominal wall; later — bladder closure, genital reconstruction, and spinal cord repair. Prognosis: survival >90\% with modern care. Bowel and bladder function outcomes are guarded; most require intermittent catheterization and bowel management.

\medterm{Colic}-- See \hyperlink{Infantile Colic}{Infantile Colic}.

\medterm{Congenital CMV}-- Cytomegalovirus infection acquired in utero, the most common congenital viral infection. Maternal primary infection poses highest risk of fetal transmission and disease. Can cause microcephaly, sensorineural hearing loss, chorioretinitis, hepatosplenomegaly, and petechiae. Diagnosis by CMV PCR in urine or saliva. Treatment: ganciclovir (IV) or valganciclovir (oral) for 6 months for symptomatic congenital CMV. Long-term sequelae include hearing loss and developmental delay.

\medterm{Congenital Diaphragmatic Hernia}-- A defect in the diaphragm allowing abdominal organs to herniate into the thoracic cavity, causing pulmonary hypoplasia and respiratory distress at birth. Most commonly a posterolateral (Bochdalek) defect on the left side. Presents with cyanosis, scaphoid abdomen, and bowel sounds in the chest. Treatment: intubation, nasogastric decompression, stabilization, and surgical repair once pulmonary hypertension is controlled. See also \hyperlink{Bochdalek Hernia}{Bochdalek Hernia} (the most common specific type).

\medterm{Congenital Heart Disease}-- Structural heart defects present at birth, the most common congenital anomaly (1\% of live births). Common defects: ventricular septal defect (VSD, 30\%), atrial septal defect (ASD), patent ductus arteriosus (PDA), tetralogy of Fallot (TOF), transposition of great arteries (TGA), coarctation of aorta, and hypoplastic left heart syndrome. Presentation depends on shunt direction and severity: cyanotic (TOF, TGA, tricuspid atresia) vs. acyanotic (VSD, ASD, PDA). Cyanotic: blue spells, clubbing, polycythemia. Acyanotic: heart murmur, heart failure, failure to thrive. Detected by prenatal ultrasound, newborn pulse oximetry screening, or clinical signs. Diagnosis: echocardiography. Treatment: surgical repair or palliation, catheter-based interventions for selected defects.

\medterm{Congenital Heart Murmur}-- A heart murmur detected in the newborn period, present in approximately 50-60\% of healthy newborns. Most are innocent (functional) murmurs that disappear by 6 months of age. Pathologic murmurs may indicate congenital heart defects. Differentiation based on murmur characteristics, timing, and associated findings. Echocardiography recommended for pathologic murmurs. Innocent murmurs require no treatment.

\medterm{Congenital Hip Dislocation}-- See \hyperlink{Hip Dysplasia}{Hip Dysplasia}.

\medterm{Congenital Pneumonia}-- Lung infection acquired in utero, often from ascending chorioamnionitis. Presents with respiratory distress, fever or hypothermia, and sepsis. Common organisms include Group B Streptococcus, E. coli, and Ureaplasma. Diagnosis by chest X-ray showing infiltrates and blood culture. Treatment: broad-spectrum IV antibiotics (ampicillin and gentamicin) plus respiratory support.

\medterm{Congenital Rubella Syndrome}-- Fetal infection with rubella virus, particularly dangerous in the first trimester. Classic triad: congenital heart disease (patent ductus arteriosus, pulmonary stenosis), cataracts, and sensorineural hearing loss. Also causes hepatosplenomegaly, thrombocytopenia, and blueberry muffin rash. Prevention by maternal vaccination before pregnancy. No specific antiviral treatment; supportive care.

\medterm{Congenital Syphilis}-- Treponema pallidum infection acquired transplacentally. Can cause stillbirth, prematurity, hepatosplenomegaly, rash, snuffles, osteochondritis, and Hutchinson teeth. Diagnosed by darkfield microscopy, VDRL/RPR in neonatal serum, or PCR. Treatment: IV penicillin G---asymptomatic congenital syphilis: 10 days; symptomatic or CNS involvement: 10-14 days. Prevention through maternal screening and treatment during pregnancy.

\medterm{Congenital Zika Syndrome}-- Fetal brain infection with Zika virus, causing microcephaly, brain abnormalities, eye defects, hearing loss, and limb contractures. Transmission occurs via mosquito bite or sexual contact. Diagnosis by Zika virus PCR in maternal serum or amniotic fluid. No specific treatment. Prevention by avoiding endemic areas during pregnancy. Surveillance and developmental monitoring recommended for exposed infants.

\medterm{Contraception for Adolescents}-- Methods to prevent unintended pregnancy. Options: combined oral contraceptives, progestin-only pills, depot medroxyprogesterone acetate, contraceptive implant (Nexplanon), intrauterine devices (IUDs), condoms, and abstinence. Long-acting reversible contraceptives (LARC) such as IUDs and implants have highest efficacy and are first-line recommendations. Dual protection against STIs requires barrier methods.

\medterm{Cornelia de Lange Syndrome}-- A genetic disorder caused by mutations in genes of the cohesin complex (NIPBL 60\%, SMC1A, SMC3, RAD21, HDAC8). Features: characteristic facies — synophrys (confluent eyebrows), long eyelashes, low-set ears, depressed nasal bridge with anteverted nares, thin lips, and micrognathia; prenatal and postnatal growth restriction; hirsutism; microcephaly; limb defects (oligodactyly, ectrodactyly, phocomelia, 30\%); and congenital heart disease (VSD, ASD). Intellectual disability ranges from moderate to severe. Gastroesophageal reflux and feeding difficulties are near-universal. Diagnosis: clinical and genetic testing. Management: growth and feeding support (gastrostomy tube often required), antireflux therapy, developmental therapies, and surveillance for hearing loss, vision problems, and scoliosis.

\medterm{Craniosynostosis}-- Premature fusion of one or more cranial sutures, restricting skull growth perpendicular to the affected suture. Presents with abnormal head shape (scaphocephaly from sagittal fusion, plagiocephaly from coronal fusion). May increase intracranial pressure. Diagnosis by physical examination and CT with 3D reconstruction. Treatment: surgical cranioplasty to release fused suture and allow normal brain growth, typically between 3-12 months of age.

\medterm{Cri-du-Chat Syndrome}-- A genetic disorder caused by a deletion of the short arm of chromosome 5 (5p-), with the size of the deletion correlating with severity. Incidence 1 in 15,000-50,000. Features: a characteristic high-pitched, cat-like cry in infancy (due to laryngeal hypoplasia), microcephaly, round face, hypertelorism, epicanthal folds, micrognathia, and hypotonia. Intellectual disability ranges from moderate to severe. Other findings: congenital heart disease (VSD, PDA), scoliosis, and hearing loss. Diagnosis: karyotype, FISH, or chromosomal microarray showing 5p deletion. Management: early intervention programs, speech therapy, physical therapy, and treatment of associated anomalies. Patients have a characteristic friendly, sociable personality. Life expectancy is near normal for mild cases.

\medterm{Croup}-- Acute laryngotracheobronchitis, most commonly caused by parainfluenza virus. Affects children 6 months to 3 years. Presents with barking cough, stridor, hoarseness, and inspiratory distress, worse at night. Diagnosis is clinical; X-ray shows subglottic narrowing (steeple sign). Treatment: single dose of oral dexamethasone and nebulized epinephrine for moderate to severe cases. Self-limited; resolves in 3-7 days.

\medterm{Crouzon Syndrome}-- An autosomal dominant craniosynostosis syndrome caused by mutations in the FGFR2 or FGFR3 genes, leading to premature fusion of multiple cranial sutures. Features: brachycephaly (short, broad head from bilateral coronal synostosis), proptosis (shallow orbits), hypertelorism, midface hypoplasia, beaked nose, and Class III malocclusion. Unlike Apert syndrome, there are no hand or foot anomalies. Intelligence is typically normal. Increased intracranial pressure can occur from progressive synostosis. Diagnosis: clinical and genetic testing; CT for craniosynostosis confirmation. Management: multidisciplinary — cranial vault remodeling surgery for craniosynostosis (3-6 months), midface advancement (Le Fort III osteotomy, 6-12 years), and orthognathic surgery for malocclusion. Lifelong monitoring for elevated ICP, sleep apnea, and vision changes.

\medterm{Cystic Hygroma}-- Congenital lymphatic malformation composed of thin-walled lymphatic cysts, manifesting as a soft, fluid-filled, painless, translucent (transilluminates brilliantly) mass, typically in the posterior triangle of the neck (75\%, left-sided), axilla (20\%), mediastinum, retroperitoneum, or groin. Most appear at birth or in early childhood; size ranges from a few cm to massive cystic masses. Histologic: dilated endothelial-lined lymphatic channels without communication to normal lymphatics. Etiology: congenital malformation of the lymphatic venous system, often from sequestration of lymphatic tissue that fails to drain into venous system, occurring during embryogenesis (~9 weeks); also reported in Turner syndrome ($\sim 50\%$ of prenatally detected cystic hygroma), Noonan syndrome, Trisomy 21, 18, 13. Diagnosis: prenatal ultrasound (nuchal translucency > 3.5 mm or septated cystic mass), fetal MRI, fetal echocardiography; karyotype and chromosomal microarray. Postnatal: physical exam, ultrasound, MRI for extent, airway assessment (cyst may compromise airway). Treatment: prenatal intervention for severe cases (sclerotherapy in utero); postnatal surgical excision (treatment of choice; complete excision if possible) preserving facial nerve, vessels, vital structures; sclerotherapy (picibanil OK-432, doxycycline, ethanol, bleomycin) for inoperable or recurrent cases; tracheostomy for airway compromise; multi-disciplinary management. Prognosis: depends on size, location, and associated anomalies; mortality high for large fetal hygroma (~80\% with hydrops); isolated small postnatal lesions surgical cure rate excellent.

\medterm{Developmental Milestones}-- Skills and abilities that children typically acquire at predictable ages. Examples: social smile at 2 months, rolling over at 4-6 months, sitting independently at 6-8 months, walking at 12 months, and speaking two-word phrases at 2 years. Delays in reaching milestones may indicate developmental delay, autism, or neurologic impairment. Screening at each well-child visit using age-appropriate tools.

\medterm{DiGeorge Syndrome}-- 22q11.2 deletion syndrome. See immunology.

\medterm{Diphtheria Vaccine}-- Toxoid vaccine containing formalin-inactivated diphtheria toxin, providing immunity against the exotoxin produced by Corynebacterium diphtheriae rather than the bacterium itself. Administered as part of combination vaccines: DTaP (pediatric, 2, 4, 6, 15-18 months, 4-6 years), Tdap (adolescent booster at 11-12 years, and during each pregnancy at 27-36 weeks), and Td (adult booster every 10 years). The toxoid is adsorbed onto an aluminium salt adjuvant to enhance immunogenicity. Antibody levels wane over 10-15 years, necessitating periodic boosters. Diphtheria toxoid is also combined with tetanus and pertussis antigens in Tdap and DTaP, and with inactivated polio vaccine in DTaP-IPV formulations. Efficacy exceeds 95 percent after a complete primary series. Before widespread vaccination, diphtheria was a leading cause of childhood mortality from respiratory obstruction and myocarditis. In India and many low- and middle-income countries, diphtheria toxoid is delivered as part of the pentavalent vaccine (DPT+HepB+Hib), using whole-cell rather than acellular pertussis, with booster doses of DPT given at 16-24 months and 5-6 years. See \hyperlink{Pentavalent Vaccine}{Pentavalent Vaccine}, \hyperlink{India Immunization Schedule}{India Immunization Schedule}, and \hyperlink{Diphtheria}{Diphtheria} in infectious disease for disease details.

\medterm{Down Syndrome}-- Trisomy 21, most common viable trisomy. See genetics.

\medterm{DTaP Vaccine}-- Diphtheria, Tetanus, and acellular Pertussis vaccine given at 2, 4, 6, and 15-18 months, and 4-6 years. Protects against three serious bacterial diseases. Diphtheria causes respiratory and cardiac complications. Tetanus causes lockjaw and muscle rigidity. Pertussis causes whooping cough, most severe in young infants. Acellular formulation reduces side effects compared to whole-cell pertussis vaccine.

\medterm{Ductal-Dependent Lesions}-- Congenital heart defects that require the ductus arteriosus to remain patent for adequate systemic or pulmonary blood flow. Includes pulmonary atresia, critical aortic stenosis, coarctation, hypoplastic left heart syndrome, and transposition of the great arteries. Present with cyanosis or shock when the ductus closes (typically within 24-48 hours). Treatment: prostaglandin E1 to maintain ductal patency until surgical intervention.

\medterm{Duodenal Atresia}-- A congenital obstruction of the duodenum caused by failure of recanalization during the 8th-10th week of gestation. Associated with Down syndrome (30\%), polyhydramnios (due to impaired fetal swallowing of amniotic fluid), and other anomalies (VACTERL, congenital heart disease, annular pancreas, malrotation). Presents within hours of birth with bilious vomiting (90\%, if distal to the ampulla of Vater), epigastric distension, and failure to pass meconium. The classic X-ray finding is the double bubble sign (air in the stomach and proximal duodenum, with no distal gas). Diagnosis: prenatal ultrasound (polyhydramnios, dilated stomach and duodenum), postnatal abdominal X-ray, and upper GI series to confirm and exclude malrotation. Treatment: nasogastric decompression, IV fluids, and surgical duodeno-duodenostomy (diamond-shaped anastomosis, Kimura technique). Prognosis: excellent with surgical repair (>95\% survival), but depends on associated anomalies.

\medterm{Dyslexia}-- A specific learning disorder with neurobiological origin, characterized by difficulties with accurate and fluent word recognition, spelling, and decoding ability, despite adequate intelligence, motivation, and educational opportunity. Affects approximately 5--10\% of the population. The underlying cause involves impaired phonological processing (difficulty recognizing and manipulating the sounds of spoken language) and deficits in rapid automatized naming. Functional neuroimaging shows reduced left temporoparietal and occipitotemporal activation during reading. Diagnosis: comprehensive educational assessment evaluating reading accuracy, fluency, comprehension, and phonological processing. Management: structured, explicit, multisensory phonics-based reading instruction (Orton-Gillingham, Wilson Reading System, Barton), assistive technology, and educational accommodations.

\medterm{Edwards Syndrome}-- Trisomy 18, the second most common autosomal trisomy in live births (1 in 5,000). Caused by complete trisomy 18 (94\%), mosaicism, or partial trisomy (translocation). Features: severe intrauterine growth restriction, clenched fists with overlapping fingers (index over middle, fifth over fourth), rocker-bottom feet, low-set malformed ears, micrognathia, prominent occiput, and congenital heart defects (VSD, PDA, polyvalvular disease, 90-95\%). Other findings: horseshoe kidney, esophageal atresia, omphalocele, and neural tube defects. Diagnosis: prenatal screening (low PAPP-A, low hCG, high inhibin A), NIPT, and karyotype/CMA. Management: comfort care is the standard (median survival 5-15 days; 5-10\% survive to 1 year). Surgical intervention for life-threatening anomalies is controversial. Survivors have severe intellectual disability.

\medterm{Ehlers-Danlos Syndrome}-- A group of hereditary connective tissue disorders caused by defects in collagen synthesis or processing. The most common types: Classic (COL5A1, COL5A2) — skin hyperextensibility, atrophic scarring, and joint hypermobility; Hypermobile (hEDS, unknown genetic cause, most common) — generalized joint hypermobility, chronic pain, and easy bruising; Vascular (COL3A1, most dangerous) — thin translucent skin, easy bruising, arterial rupture (carotid, aortic, visceral), uterine rupture, and sigmoid colon perforation; Kyphoscoliotic (PLOD1) — congenital hypotonia, scoliosis, and joint laxity; and Arthrochalasia (COL1A1, COL1A2) — severe joint hypermobility and congenital hip dislocation. Diagnosis: clinical (Beighton score for hypermobility), collagen typing on skin biopsy (vascular type), and genetic testing. Management: activity modification and joint protection, vascular surveillance for vascular EDS (carotid duplex, aortic imaging, avoidance of invasive procedures), high-dose vitamin C, and genetic counseling. No cure; management is symptomatic and preventive.

\medterm{Emergency Contraception}-- Postcoital contraception to prevent pregnancy after unprotected intercourse or contraceptive failure. Levonorgestrel (Plan B): effective within 72 hours, available over the counter. Ulipristal acetate (Ella): effective within 120 hours, prescription only. Copper IUD: most effective emergency contraception, placed within 5 days. Not abortifacients; prevent or delay ovulation.

\medterm{Enuresis (Bedwetting)}-- Involuntary urination during sleep in a child aged 5 years or older, affecting approximately 5--10\% of 7-year-olds. Primary enuresis: the child has never achieved sustained nighttime dryness. Secondary enuresis: wetting recurs after at least 6 months of dryness, often triggered by psychosocial stress, urinary tract infection, or diabetes. Pathophysiology involves a mismatch between nocturnal bladder capacity, nighttime urine production (lack of circadian ADH rhythm), and inability to arouse from sleep in response to bladder fullness. Evaluation: history (wetting pattern, fluid intake, constipation), urinalysis (glucose, infection). Management: behavioral (timed voiding, reward systems, fluid restriction before bed), enuresis alarm (first-line, ~70\% success), and pharmacotherapy (desmopressin--synthetic ADH--reduces nighttime urine production).

\medterm{Epiglottitis}-- Inflammation of the epiglottis, a life-threatening emergency. Most commonly caused by Haemophilus influenzae type b (now rare due to vaccination). Presents with high fever, drooling, dysphagia, tripod position, and rapid progression to airway obstruction. Do not examine the throat. Treatment: secure airway in operating room, IV antibiotics (ceftriaxone). Requires immediate airway management; mortality high if untreated.

\medterm{Epispadias}-- A congenital anomaly where the urethral meatus opens on the dorsal (top) surface of the penis in boys or in a split, open position on the clitoral hood in girls. Caused by abnormal development of the genital tubercle and cloacal membrane. In boys, the opening may be at the glans (glandular, mildest), penile shaft, or penopubic (most severe, associated with bladder exstrophy). In girls, the urethra is short with a bifid clitoris and separated labia. Epispadias presents as part of the exstrophy-epispadias complex spectrum. The penis is short, wide, and upwardly curved (dorsal chordee). Diagnosis: clinical. Treatment: surgical repair (Mitchell or Cantwell-Ransley technique) between 6-12 months — urethral reconstruction, penile straightening, and urethroplasty. Prognosis: urinary continence is achieved in 70-80\% with repaired epispadias.

\medterm{Erb's Palsy}-- Upper brachial plexus injury (C5-C6) causing weakness or paralysis of the shoulder and arm muscles. Presents with the characteristic waiter's tip position: arm adducted, elbow extended, forearm pronated, wrist flexed. Most commonly caused by shoulder dystocia during delivery. See also \hyperlink{Brachial Plexus Injury}{Brachial Plexus Injury}.

\medterm{Esophageal Atresia}-- A congenital anomaly where the esophagus ends in a blind pouch, often with a tracheoesophageal fistula connecting the distal esophagus to the trachea. Presents with drooling, inability to pass a nasogastric tube, and choking with feedings. The 3 Es association: Esophageal atresia, Esophageal anomalies, and Extremity anomalies. Treatment: surgical ligation of fistula and esophageal anastomosis.

\medterm{Exchange Transfusion}-- Emergent procedure to rapidly lower severe hyperbilirubinemia by removing small volumes of the infant's blood and replacing with compatible donor blood. Used when phototherapy fails or bilirubin approaches toxic levels (greater than 25 mg/dL in term infants or above exchange threshold on Bhutani curve). Risks include infection, electrolyte imbalances, and air embolism. Double-volume exchange (160-200 mL/kg) replaces approximately 85\% of circulating red blood cells.

\medterm{Failure to Thrive}-- Inadequate growth in children, defined as weight below the 5th percentile or crossing two major percentile lines on growth charts. Causes include caloric insufficiency (neglect, poor feeding), organic disease (congenital heart disease, renal disease), and psychosocial factors. Evaluation includes comprehensive history, physical examination, laboratory studies, and calorie counts. Treatment addresses underlying cause and nutritional rehabilitation.

\medterm{Fanconi Anemia}-- An autosomal recessive (or X-linked) bone marrow failure syndrome caused by mutations in at least 22 FANC genes involved in DNA repair (the FA/BRCA pathway). Features: progressive pancytopenia (onset typically at 5-10 years), radial ray anomalies (thumb hypoplasia or aplasia, radial club hand), microcephaly, short stature, café-au-lait spots, and renal anomalies. Increased risk of myelodysplastic syndrome, acute myeloid leukemia, and solid tumors (head and neck squamous cell carcinoma, gynecological cancer). Diagnosis: chromosomal breakage test (diepoxybutane or mitomycin C — increased chromosome breaks) and FANC gene sequencing. Management: androgens (oxymetholone) for cytopenias, hematopoietic stem cell transplant for bone marrow failure (with reduced-intensity conditioning), and lifelong cancer surveillance.

\medterm{Febrile Seizures}-- Seizures occurring in children 6 months to 5 years triggered by fever (temperature greater than 38 degrees Celsius), without intracranial infection or defined cause. Occurs in 2-5\% of children. Simple: generalized, single seizure lasting less than 15 minutes, no recurrence within 24 hours. Complex: focal, prolonged (greater than 15 minutes), or recurrent within 24 hours. Diagnosis of exclusion after ruling out meningitis and intracranial infection. In simple febrile seizures, there is no increased risk of epilepsy. Management: parental education on seizure first aid, antipyretics for comfort (do not prevent recurrence), and reassurance. No routine antiepileptic medication is indicated. Lumbar puncture should be considered when meningeal signs are present or in infants under 12 months with incomplete vaccination.

\medterm{Fetal Distress}-- Abnormal fetal status during labor, indicated by abnormal fetal heart rate patterns (late or prolonged decelerations, minimal variability, absent accelerations). May indicate fetal hypoxia, cord compression, or placental insufficiency. Management includes intrauterine resuscitation: change maternal position, stop oxytocin, give IV fluids and oxygen, and amnioinfusion. If persistent, emergent delivery by cesarean or operative vaginal delivery.

\medterm{Fifth Disease}-- Erythema infectiosum caused by Parvovirus B19. Characterized by bright red cheeks (slapped cheek appearance) followed by a lacy, reticulated rash on the trunk and extremities. Mild febrile illness in children. Concern in pregnancy (hydrops fetalis risk) and immunocompromised patients (pure red cell aplasia). Diagnosis by serology. Self-limited; no specific treatment.

\medterm{Fine Motor Skills}-- The ability to use small muscles (hands/fingers) for precise movements, developing with age: reaching (3-4 months), grasping (5-6 months), pincer grasp (9-12 months), self-feeding (12-18 months), building towers/scribbling (18-24 months), drawing shapes and using utensils (3-5 years), and handwriting. Delays relate to neuromotor conditions (cerebral palsy, developmental delay), visual problems, and autism (with other signs). Assessed at well-child visits; early referral for occupational therapy improves function. Monitor bimanual use, hand dominance, and age-appropriate play skills.

\medterm{Fontanelle}-- Soft membranous spaces between skull bones that allow brain growth and skull molding during delivery. Anterior fontanelle diamond-shaped, closes by 12-18 months. Posterior fontanelle triangular, closes by 2-3 months. Bulging fontanelle suggests increased intracranial pressure; sunken fontanelle indicates dehydration. Palpation is part of the routine infant physical examination.

\medterm{Formula Feeding}-- Alternative to breastfeeding using commercially prepared infant formula. Types include cow's milk-based, soy-based, hydrolyzed protein, and amino acid-based formulas. Iron-fortified formulas standard to prevent iron deficiency anemia. Preparation requires clean water and proper sterilization of bottles. Feeding on demand approximately every 3-4 hours. Provides adequate nutrition when breastfeeding is not possible or desired.

\medterm{Fragile X Syndrome}-- X-linked dominant disorder caused by CGG trinucleotide repeat expansion in the FMR1 gene. Most common inherited cause of intellectual disability. Features: long face, prominent ears, macroorchidism in post-pubertal males, and behavioral problems including autism spectrum features. Diagnosis by DNA testing for CGG repeat expansion. No cure; behavioral and educational interventions.

\medterm{Functional Asplenia}-- Loss of splenic function due to sickling and infarction in sickle cell disease, typically by age 5 years. Increases risk of overwhelming bacterial infection, particularly from encapsulated organisms (Streptococcus pneumoniae, Haemophilus influenzae). Management: prophylactic penicillin, pneumococcal vaccines, and parental education about fever as a medical emergency.

\medterm{Galactocele}-- See obstetrics and gynecology.

\medterm{Galant Reflex}-- A primitive reflex present from birth to 4 months, elicited by stroking the skin along the side of the infant's back. The infant curves toward the stimulated side. Assessed by running a finger from the lateral back toward the hip while the infant is in ventral suspension. Absence may indicate neurologic abnormality. One of the trunk reflexes evaluated during newborn neurologic examination.

\medterm{Gastroschisis}-- A full-thickness abdominal wall defect, usually to the right of the umbilicus, through which bowel herniates without a covering membrane. Not associated with chromosomal abnormalities. Presents at birth with exposed bowel loops that may be edematous, thickened, and discolored from amniotic fluid exposure. Management: immediate coverage with saline-moistened dressings, nasogastric decompression, and surgical repair (primary closure or staged reduction).

\medterm{German Measles (Rubella)}-- See infectious disease.

\medterm{Glandular Fever (Infectious Mononucleosis)}-- See infectious disease.

\medterm{Glycogen Storage Disease}-- A group of inherited metabolic disorders caused by defects in glycogen synthesis, breakdown, or regulation, resulting in abnormal glycogen accumulation in tissues. Major types: Type I (von Gierke, G6PC deficiency) — fasting hypoglycemia, lactic acidosis, hepatomegaly, hyperuricemia, and hypertriglyceridemia. Type II (Pompe disease, GAA deficiency) — infantile form: hypertrophic cardiomyopathy, hypotonia, and death by 1 year; late-onset: proximal myopathy and respiratory weakness. Type III (Cori disease, AGL deficiency) — hypoglycemia, hepatomegaly, and myopathy. Type IV (Andersen disease, GBE1 deficiency) — liver failure, cirrhosis, and cardiomyopathy. Type V (McArdle disease, PYGM deficiency) — exercise intolerance, myoglobinuria, and second-wind phenomenon. Type VI (Hers disease, PYGL deficiency) — mild fasting hypoglycemia and hepatomegaly. Type VII (Tarui disease, PFKM deficiency) — exercise intolerance similar to McArdle with hemolytic anemia. Diagnosis depends on the type: enzyme assay in liver or muscle biopsy, genetic testing, and newborn screening for Pompe. Treatment varies by type: Type I — continuous glucose therapy (cornstarch), allopurinol, and lipid-lowering agents. Type II (Pompe) — enzyme replacement therapy (alglucosidase alfa). Types III, VI — high-protein diet. Types V, VII — exercise modification and creatine supplementation.

\medterm{Goldenhar Syndrome}-- A variant of hemifacial microsomia characterized by the triad: epibulbar dermoids (benign conjunctival tumors), preauricular skin tags or sinuses, and vertebral anomalies (hemivertebrae, fused vertebrae, scoliosis). Also involves facial asymmetry, microtia, conductive hearing loss, and macrostomia. Other associations: congenital heart disease (VSD, tetralogy of Fallot), renal anomalies, and cleft lip/palate. Most cases are sporadic. Diagnosis: clinical and imaging (CT for skeletal anomalies, echocardiogram, renal ultrasound). Management: multidisciplinary — ophthalmology (epibulbar dermoid excision), audiology (hearing aids), plastic surgery (auricular reconstruction, mandibular distraction), and orthopedic monitoring for scoliosis.

\medterm{Gonorrhea}-- See infectious disease.

\medterm{Gross Motor Skills}-- The ability to use large muscle groups for movement and posture, developing in a predictable sequence in infancy/childhood: head control (3 months), sitting (6 months), crawling (9 months), standing (10-12 months), walking (12-18 months), running/jumping (2-3 years), hopping and balance (3-5 years). Delays warrant assessment: cerebral palsy, muscular dystrophy, developmental delay, hypotonia, and environmental factors. Milestone monitoring at well-child visits (screening—ASQ) and early intervention improve outcomes; gross motor regression is always concerning. Encourage tummy time and active play.

\medterm{Group B Strep}-- Group B Streptococcus (GBS, Streptococcus agalactiae): a gram-positive coccus colonizing the GI and genitourinary tract in 10-30 percent of women; major cause of neonatal sepsis and meningitis. Neonatal disease: early-onset (<7 days—vertical transmission during delivery: sepsis, pneumonia, meningitis) and late-onset (7 days-3 months—sepsis, meningitis, osteomyelitis). Risk: maternal GBS carriage, preterm, prolonged rupture of membranes, intrapartum fever, and previous GBS-affected infant. Prevention: universal screening (35-37 weeks) and intrapartum antibiotic prophylaxis (penicillin) for carriers. Diagnosis: blood/CSF culture, PCR. Management: empiric antibiotics (penicillin/ampicillin + gentamicin), supportive care; also causes UTI and chorioamnionitis in pregnancy.

\medterm{Growing Pains}-- Recurrent, bilateral leg pain in children ages 3-12, typically occurring in the late afternoon or evening and resolving by morning. Pain is deep, aching in the calves, shins, and thighs (not joints). No limping, fever, swelling, or functional impairment. Etiology unclear: may relate to increased activity, musculoskeletal overuse, or decreased pain threshold. Diagnosis: clinical (diagnosis of exclusion). Red flags suggesting pathology: unilateral pain, joint involvement, morning stiffness, limping, fever, night sweats, nocturnal pain waking the child. Treatment: reassurance, massage, stretching, warm baths, and acetaminophen/ibuprofen as needed. Benign and self-limited.

\medterm{Growth Chart}-- A graphical representation of growth (weight, height/length, head circumference, BMI) plotted against age, standardized by population (WHO growth standards 0-5 years; CDC growth charts 2-20 years; percentiles/z-scores). Used to monitor normal growth, detect growth disorders (failure to thrive, short stature, obesity, microcephaly/macrocephaly), and assess nutritional status. Interpretation: track crossing of percentiles (crossing downward—poor growth; upward—overfeeding/obesity), symmetry, and parental heights (target height). Regular plotting at well-child visits identifies deviations early, prompting evaluation (nutrition, endocrine, chronic disease, genetic).

\medterm{Gynecomastia}-- Enlargement of breast tissue in males, often bilateral, occurring physiologically in neonates, during puberty (40-60\% of boys), and in aging men. Pathologic causes include Klinefelter syndrome, testicular tumors, medications (spironolactone, cannabis), and liver disease. Persistent or asymmetric gynecomastia warrants evaluation. Management: observation for physiologic cases; treatment of underlying cause; mastectomy for severe persistent cases.

\medterm{Hand Foot and Mouth Disease}-- A common viral illness (mostly children <10 years) caused by coxsackievirus A16 (and enterovirus 71), transmitted by fecal-oral/respiratory. Features: fever, sore throat, and a vesicular rash on the hands, feet, and mouth (ulcers on buccal mucosa, tongue, palate; vesicles on palms/soles); enterovirus 71 may cause severe CNS disease (meningoencephalitis, polio-like paralysis) and cardiorespiratory failure. Diagnosis: clinical. Management: supportive—hydration, analgesia, and antipyretics; watch for dehydration; severe/neurologic cases (EV71) require hospital care. Highly contagious; hygiene and isolation. Usually self-limiting in 7-10 days.



\medterm{Head Circumference}-- Measured at the largest circumference above the ears and eyebrows, reflecting brain growth. Plotted on growth charts to detect microcephaly (below 3rd percentile) or macrocephaly (above 98th percentile). Rapid increase may indicate hydrocephalus or intracranial hemorrhage; failure to grow may indicate microcephaly or brain atrophy. Measured at each well-child visit until age 2.

\medterm{Head lice}-- An infestation of the scalp by the obligate human ectoparasite Pediculus humanus capitis, a wingless, blood-sucking insect that feeds on human blood, most common in children aged 3-12. It is extremely common, especially in school-aged children aged 3-11 years, with an estimated 6-12 million infestations annually in the United States. Transmission is primarily through direct head-to-head contact (the most common route), and less commonly through shared combs, brushes, hats, headphones, pillows, or other personal items. Lice do not jump or fly; they crawl. Diagnosis is confirmed by finding live lice or nits attached to hair shafts within 1 cm of the scalp using a fine-toothed detection comb. Treatment: topical pediculicides containing permethrin 1\% or dimeticone; oral ivermectin for treatment-resistant cases; wet combing (bug busting) as a non-insecticidal alternative. Bedding and clothing should be washed in hot water (60 degrees Celsius) and non-washable items sealed in a plastic bag for 2 weeks. Prophylactic treatment of household contacts is not recommended. School exclusion policies vary; the American Academy of Pediatrics and the UK National Institute for Health and Care Excellence recommend children remain in school after first treatment, as no-nit policies are unnecessary and promote stigma.

\medterm{Hearing Screening}-- Newborn hearing screening (universal): otoacoustic emissions (OAE) and/or automated auditory brainstem response (AABR) before discharge; referral for comprehensive audiology if failed. Early detection (<6 months) of congenital hearing loss (1-3/1000) allows early intervention (hearing aids, cochlear implants, early communication) with vastly improved language outcomes. Risk factors: family history, prematurity, low birth weight, infections (CMV, meningitis), ototoxic drugs, and jaundice. Ongoing surveillance: behavioral screening and audiometry at well-child visits and school; signs of delayed speech/response to sound prompt referral. Follow-up and documentation are critical.

\medterm{Hemifacial Microsomia}-- A congenital craniofacial disorder involving asymmetric hypoplasia of the structures derived from the first and second branchial arches, most commonly affecting the mandible, ear, and soft tissues on one side. Also known as craniofacial microsomia or Goldenhar syndrome when epibulbar dermoids and vertebral anomalies are present. Presents with asymmetry: mandibular hypoplasia (Pruzansky-Kaban classification I-III), microtia (small or absent pinna), preauricular tags or sinuses, conductive hearing loss, and macrostomia (enlarged oral commissure). Severity ranges from mild asymmetry to severe facial deformity. Diagnosis: clinical, CT for skeletal assessment, and audiometry. Management: reconstructive surgery (mandibular distraction, costochondral graft, or orthognathic surgery), auricular reconstruction, hearing rehabilitation, and speech therapy.

\medterm{Henoch-Schonlein Purpura}-- See \hyperlink{IgA Vasculitis (Henoch-Schonlein Purpura)}{IgA Vasculitis}.

\medterm{Hepatitis B Vaccine}-- Recombinant vaccine given at birth, 1-2 months, and 6-18 months. Protects against hepatitis B virus, which can cause chronic liver disease, cirrhosis, and hepatocellular carcinoma. Also prevents perinatal transmission when given to infants of HBsAg-positive mothers within 12 hours of birth, plus hepatitis B immune globulin.

\medterm{Hib}-- Haemophilus influenzae type b: an encapsulated gram-negative bacterium causing serious invasive disease in unvaccinated children—meningitis, epiglottitis, pneumonia, septic arthritis, cellulitis, and sepsis. Highly prevented by the Hib conjugate vaccine (given in infancy, part of routine schedule). Invasive Hib disease in vaccinated populations is rare. Diagnosis: culture/PCR (blood, CSF). Management: antibiotics (third-generation cephalosporin—ceftriaxone), and respiratory/droplet precautions; treatment of close contacts with rifampin chemoprophylaxis (household/school, especially with young children); monitor for complications (meningitis—deafness). Epiglottitis is an airway emergency.

\medterm{Hip Dysplasia}-- Developmental dysplasia of the hip (DDH): abnormal development of the acetabulum/femoral head causing instability, subluxation, or dislocation. Risk factors: female, breech presentation, family history, firstborn, and tight swaddling. Features: asymmetric skin creases, limited abduction, Ortolani (relocation) and Barlow (dislocation) tests in the newborn; later—gait limp, leg length discrepancy, Trendelenburg. Diagnosis: ultrasound (age <6 months), X-ray (older). Management: early—Pavlik harness (success >90 percent); late—closed reduction/casting or surgery. Screening (clinical ± selective ultrasound) and early referral prevent long-term arthritis and disability.

\medterm{Holt-Oram Syndrome}-- An autosomal dominant disorder caused by TBX5 gene mutations, characterized by congenital heart disease and upper limb malformations. Cardiac anomalies: atrial septal defect (most common), VSD, and conduction abnormalities (heart block, atrial fibrillation). Limb defects are bilateral and asymmetric: thumb anomalies (absent, hypoplastic, triphalangeal, or finger-like thumb, 100\%), radial dysplasia (radial club hand), and phocomelia. Diagnosis: clinical (heart defect + preaxial radial ray anomaly) and TBX5 genetic testing. Management: cardiac repair (ASD/VSD closure, pacemaker for heart block), reconstructive hand surgery (pollicization for absent thumb, centralization for radial club hand), and genetic counseling.

\medterm{Homocystinuria}-- An autosomal recessive inborn error of methionine metabolism caused most commonly by deficiency of cystathionine beta-synthase (CBS), leading to accumulation of homocysteine and methionine. Features (develop in childhood if untreated): ectopia lentis (lens dislocation, typically downward, differentiating from Marfan syndrome), severe myopia, intellectual disability, thromboembolism (the major cause of morbidity/mortality, affecting arteries and veins), osteoporosis, and a Marfanoid habitus (tall stature, pectus deformities, arachnodactyly). Diagnosis: elevated plasma homocysteine and methionine, homocystinuria, and CBS gene sequencing. Newborn screening includes homocystinuria in expanded panels. Treatment: pyridoxine (vitamin B6) responsiveness testing (approximately 50\% are responsive), betaine (to remethylate homocysteine to methionine), methionine-restricted diet (especially for non-responders), and anticoagulation (aspirin for thromboprophylaxis). Lifelong monitoring for thromboembolism and lens dislocation.

\medterm{Horseshoe Kidney}-- The most common renal fusion anomaly (1 in 400-500), where the lower poles of both kidneys are fused across the midline by an isthmus of renal parenchyma or fibrous tissue. The fused kidney is located lower in the abdomen than normal (at L3-L5, anterior to the aorta and IVC) due to the isthmus being trapped under the inferior mesenteric artery during ascent. Associated with Turner syndrome, Edwards syndrome, Patau syndrome, and other urogenital anomalies. Often asymptomatic and discovered incidentally. May present (30\%) with hydronephrosis from UPJ obstruction (the ureter passes anterior to the isthmus, causing kinking), recurrent UTIs, nephrolithiasis, and increased risk of Wilms tumor. Diagnosis: ultrasound, CT, or IV pyelogram. Treatment: surgical management of complications (pyeloplasty for UPJ obstruction, and nephrolithotomy for stones). Surveillance imaging for Wilms tumor is controversial.

\medterm{HPV Vaccine}-- Human Papillomavirus vaccine given at 11-12 years (can start at 9), with series completion by 15 years. Two-dose series if started before 15; three-dose series if started at 15 or older. Protects against HPV types causing cervical cancer, oropharyngeal cancer, and genital warts. Gardasil 9 covers nine HPV types. Expected to prevent over 90\% of HPV-associated cancers.

\medterm{Human Papillomavirus}-- The most common STI, with over 100 types. High-risk types (16, 18) cause cervical, oropharyngeal, anal, and genital cancers. Low-risk types (6, 11) cause genital warts. Prevention: HPV vaccine (Gardasil 9) protects against 9 HPV types. Screening: Pap smear and HPV testing for cervical cancer prevention. Most infections clear spontaneously.

\medterm{Inactivated Polio Vaccine (IPV)}-- Salk vaccine, a formalin-inactivated preparation containing all three poliovirus serotypes (1, 2, and 3), administered by intramuscular injection. The only polio vaccine used in the United States since 2000, and the recommended vaccine for immunocompromised individuals and pregnant women worldwide. The current CDC schedule specifies doses at 2 months, 4 months, 6-18 months, and a booster at 4-6 years. Unlike oral polio vaccine (OPV), IPV does not carry the risk of vaccine-associated paralytic polio (VAPP) or circulating vaccine-derived poliovirus (cVDPV), because the inactivated virus cannot replicate in the gut. However, IPV generates weaker intestinal mucosal immunity than OPV, meaning vaccinated individuals can still harbour and transmit wild poliovirus in the gut. IPV is included in combination vaccines such as DTaP-IPV and DTaP-IPV-Hib. Post-2016, the type 2 component was removed from both IPV (now bIPV) and OPV in response to the global eradication of wild poliovirus type 2. Efficacy exceeds 99 percent after a complete primary series. In India, IPV is given as fractional doses (fIPV) rather than full intramuscular doses. See \hyperlink{Fractional IPV (fIPV)}{Fractional IPV (fIPV)}, \hyperlink{Polio}{Polio}, and \hyperlink{Poliomyelitis}{Poliomyelitis} in infectious disease.

\medterm{Fractional IPV (fIPV)}-- A reduced-dose schedule of inactivated polio vaccine used in the WHO and India Universal Immunization Schedule, where one-fifth of the standard intramuscular IPV dose (0.1 mL instead of 0.5 mL) is given by intradermal injection at each of two deltoid sites (total 0.2 mL per fractional dose), rather than the full 0.5 mL intramuscular dose. The schedule comprises two fractional doses at 6 weeks and 14 weeks, with a third full intramuscular IPV dose at 9-12 months. This approach was adopted to address global IPV supply constraints while maintaining adequate seroconversion rates (greater than 90 percent for all three serotypes after three doses). Intradermal delivery targets the antigen-rich dermal layer, enhancing immune stimulation per microgram of antigen. Seroconversion rates with fractional IPV are comparable to full-dose IPV after three doses, though antibody titres may be slightly lower after two fractional doses alone. India adopted fIPV in 2016 as part of its dual polio vaccination strategy combining OPV and IPV. The approach is endorsed by the WHO Strategic Advisory Group of Experts on Immunization (SAGE).

\medterm{Hunter Syndrome}-- Mucopolysaccharidosis type II (MPS II), an X-linked recessive lysosomal storage disorder caused by iduronate-2-sulfatase deficiency (IDS gene mutations). Accumulation of dermatan sulfate and heparan sulfate (same as Hurler). Two forms: severe (neuronopathic, 2/3 of cases) — progressive intellectual disability, neurodegeneration, and death by 10-15 years; attenuated (no or minimal neurologic involvement) — survival into adulthood. Features: coarse facies, corneal clouding (less prominent than Hurler), hepatosplenomegaly, recurrent hernias, dysostosis multiplex, joint stiffness, short stature, and cardiac valvular disease. A characteristic feature is the pebbled, ivory-white skin lesions (papules) over the back and shoulders (pathognomonic). Diagnosis: urine GAG analysis, iduronate-2-sulfatase enzyme assay, and IDS gene sequencing. Treatment: enzyme replacement therapy (idursulfase, Elaprase) improves somatic features, and intrathecal idursulfase is under investigation for neurologic disease. HSCT is less effective than in MPS I.

\medterm{Japanese Encephalitis (JE) Vaccine}-- An inactivated (SA 14-14-2 strain, widely used) or live attenuated vaccine protecting against Japanese encephalitis, the most common cause of viral encephalitis in Asia, transmitted by Culex mosquitoes in rural rice-growing and pig-farming areas. Most infections are subclinical or cause mild febrile illness, but symptomatic disease carries a case fatality rate of 20-30 percent and neurological sequelae in 30-50 percent of survivors. In India's Universal Immunization Schedule, JE vaccine is given in districts with high JE burden: the first dose at 9-12 months and a second dose at 16-24 months, administered by subcutaneous injection. The vaccine is also part of pre-travel recommendations for long-term travellers to endemic areas (multiple doses required). Adverse effects are generally mild: injection site pain, fever, and headache. The WHO recommends JE vaccination for all children living in or travelling to endemic areas.

\medterm{Hurler Syndrome}-- Mucopolysaccharidosis type I (MPS I), the severe form, an autosomal recessive lysosomal storage disorder caused by alpha-L-iduronidase deficiency (IDUA gene mutations). Accumulation of dermatan sulfate and heparan sulfate leads to progressive multisystem involvement. Features: coarse facies (gargoyle-like appearance), corneal clouding, hepatosplenomegaly, recurrent hernias, dysostosis multiplex (kyphoscoliosis, J-shaped sella, paddle ribs, claw-hand deformity, hip dysplasia), aortic valve disease, and severe intellectual disability. Hearing loss, obstructive airway disease, and hydrocephalus (communicating) are common. Diagnosis: urine glycosaminoglycan (GAG) analysis (elevated dermatan and heparan sulfate) and IDUA enzyme assay in leukocytes. Newborn screening is available. Treatment: hematopoietic stem cell transplant (HSCT) performed early (<2.5 years) preserves cognitive function and improves somatic features but does not fully correct skeletal or corneal disease. Enzyme replacement therapy (laronidase, iduronidase) improves somatic features but does not cross the blood-brain barrier. Prognosis: untreated, death by 6-10 years from cardiorespiratory failure.

\medterm{Hydrocephalus}-- Abnormal accumulation of cerebrospinal fluid (CSF) within the ventricles of the brain, leading to increased intracranial pressure. See neurology.

\medterm{Hypertelorism}-- Increased distance between the orbits (the bony eye sockets), measured as an increased interorbital distance on imaging. Can result from premature fusion of the cranial sutures (craniosynostosis syndromes), frontonasal dysplasia, facial clefting, or can be associated with genetic syndromes (Apert, Crouzon, Noonan, 22q11.2 deletion, and CHARGE). Presents with widely spaced eyes (telecanthus). Diagnosis: clinical measurement of the intercanthal distance and orbital imaging (CT or MRI) to confirm increased interorbital distance. Management: surgical correction by orbital box osteotomy or facial bipartition for severe cases; treatment of the underlying syndrome. Not typically a functional concern.

\medterm{IgA Vasculitis (Henoch-Schonlein Purpura)}-- The most common systemic vasculitis in children (peak age 2-8 years), characterized by IgA1-containing immune complex deposition in small vessels. Often preceded by an upper respiratory tract infection (e.g., group A Streptococcus). Classic clinical tetrad: 1) Palpable purpura (non-thrombocytopenic) on lower extremities and buttocks, 2) Arthralgia/arthritis (knees, ankles), 3) Abdominal pain (colicky, risk of intussusception), and 4) Renal involvement (hematuria, proteinuria, IgA nephropathy). Lab workup: normal platelet count, elevated serum IgA, normal complement levels. Management: supportive care, NSAIDs for joint pain; systemic corticosteroids for severe abdominal pain or crescentic glomerulonephritis. Excellent prognosis with >90\% spontaneous resolution.

\medterm{Immune Thrombocytopenic Purpura}-- Autoimmune destruction of platelets by antiplatelet antibodies, causing isolated thrombocytopenia. Most common in children 2-6 years, often following viral infection. Presents with petechiae, purpura, and mucosal bleeding. Diagnosis of exclusion. Treatment: observation for mild cases; IVIG or corticosteroids for severe thrombocytopenia or active bleeding. Spontaneous resolution in 80\% within 6 months.

\medterm{Imperforate Anus}-- A congenital anomaly where the anus is completely absent or fails to open at the perineum. Classified as low (distal to the levator ani complex) or high (proximal to the levator ani). Often associated with other anomalies including vertebral, anal, cardiac, tracheoesophageal, renal, and limb defects. Presents with absence of anal opening on physical examination. Management: colostomy for high lesions; perineal anoplasty for low lesions.

\medterm{Impetigo}-- A superficial bacterial skin infection, most commonly caused by Staphylococcus aureus or Group A Streptococcus. Presents with honey-colored crusted lesions (nonbullous) or fluid-filled blisters (bullous). Common in children. Diagnosis is clinical. Treatment: topical mupirocin for localized disease; oral antibiotics for extensive disease. Prevention by good hand hygiene and avoiding sharing personal items.

\medterm{Infantile Colic}-- A behavioral syndrome in healthy infants characterized by episodes of excessive, inconsolable crying for >3 hours per day, >3 days per week, for >3 weeks (Wessel criteria). Onset typically at 2-3 weeks, peaks at 6-8 weeks, and resolves by 3-4 months of age. Etiology poorly understood: possible GI (gas, gut microbiota, cow milk protein intolerance), neurodevelopmental, or behavioral. Diagnosis: clinical; must rule out organic causes (intussusception, GERD, infection, hair tourniquet, corneal abrasion, testicular torsion). Treatment: parental reassurance, carrying/rocking, white noise, and swaddling. Probiotics, simethicone, and hypoallergenic formula may help some infants. Prognosis: excellent, self-limited.

\medterm{Influenza Vaccine}-- Inactivated influenza vaccine given annually starting at 6 months of age. Two doses needed in the first season for children under 9. Live attenuated nasal spray vaccine (FluMist) available for children over 2. Updated annually to match circulating strains. Reduces influenza complications, hospitalizations, and deaths. Recommended for all children.

\medterm{Intestinal Atresia}-- A congenital obstruction of the small intestine (jejunum or ileum) caused by an intrauterine vascular accident (mesenteric ischemia) leading to aseptic necrosis and resorption of a segment of bowel. In contrast to duodenal atresia, it is not associated with Down syndrome. Presents in the first 24-48 hours with bilious vomiting and progressive abdominal distension. Failure to pass meconium is universal with distal atresia. Classification (Grosfeld): Type I — mucosal web with intact bowel wall; Type II — blind ends connected by a fibrous cord; Type IIIa — complete separation of blind ends with a V-shaped mesenteric defect; Type IIIb — apple-peel or Christmas-tree deformity (extensive atresia with a spiral-shaped distal bowel); Type IV — multiple atresias. Diagnosis: abdominal X-ray shows dilated proximal bowel loops (number indicates level of atresia) with no distal gas; contrast enema shows an unused microcolon. Treatment: primary anastomosis after resection of the dilated proximal segment; tapering enteroplasty may be needed for severe proximal dilation. Prognosis: excellent for Type I-IIIa; guarded for IIIb and IV.

\medterm{Intraventricular Hemorrhage}-- Bleeding into the germinal matrix and ventricular system, primarily in premature infants less than 32 weeks gestation. Graded I-IV by severity on cranial ultrasound. Grade I (germinal matrix only) and Grade II (intraventricular without dilation) have good prognosis. Grade III (with ventricular dilation) and Grade IV (parenchymal involvement) risk neurodevelopmental impairment. Prevention: antenatal corticosteroids, avoiding rapid fluid shifts.

\medterm{Jaundice in newborns}-- See \hyperlink{Neonatal Jaundice (Jaundice of the Newborn)}{Neonatal Jaundice (Jaundice of the Newborn)}.

\medterm{Juvenile Idiopathic Arthritis (JIA)}-- A heterogeneous group of autoimmune and inflammatory conditions causing persistent arthritis (joint swelling, pain, limited range of motion) in children aged <16 years, lasting >6 weeks. Classification: oligoarticular (most common, 50\%, <5 joints, high risk of uveitis), polyarticular (RF-positive or RF-negative, 30\%), systemic-onset (Still disease, high spiking fevers, salmon-pink evanescent rash, serositis, hepatosplenomegaly), enthesitis-related (enthesitis, sacroiliitis, HLA-B27 association), and psoriatic arthritis. etiology: complex genetic predisposition (HLA class II associations) plus environmental triggers. Management: NSAIDs (first-line), intra-articular corticosteroid injections, methotrexate (first-line DMARD), biologic agents (anti-TNF--etanercept, adalimumab; anti-IL-6--tocilizumab for systemic JIA; anti-IL-1--anakinra). Uveitis screening is essential (slit-lamp every 3-4 months) as most cases are asymptomatic and can cause blindness.

\medterm{Kawasaki Disease}-- An acute vasculitis of childhood affecting medium-sized arteries, particularly coronary arteries. Diagnosis requires fever greater than 5 days plus 4 of 5 criteria: bilateral conjunctivitis, oral mucosal changes, cervical lymphadenopathy, extremity changes, and polymorphous rash. Complications include coronary artery aneurysms. Treatment: IV immunoglobulin (IVIG) 2 g/kg as a single infusion within the first 10 days of illness, plus aspirin---high-dose (80-100 mg/kg/day in 4 divided doses) during the acute febrile phase, then low-dose (3-5 mg/kg/day for 6-8 weeks, or longer if coronary aneurysm present). Coronary artery aneurysms occur in 15-25 percent of untreated cases; IVIG reduces this to less than 5 percent. Most common acquired heart disease in children in developed countries.

\medterm{Kernicterus}-- Bilirubin encephalopathy resulting from severe unconjugated hyperbilirubinemia (typically greater than 25 mg/dL in term infants). Bilirubin deposits in the basal ganglia, hippocampus, and brainstem nuclei. Presents with lethargy, poor feeding, high-pitched cry, opisthotonus, seizures, and choreoathetosis. May cause permanent neurologic damage including hearing loss, cerebral palsy, and intellectual disability. Prevention by treating severe jaundice with phototherapy or exchange transfusion.

\medterm{Klippel-Trenaunay Syndrome}-- A congenital vascular disorder characterized by the triad: capillary malformations (port-wine stains), venous malformations (varicose veins, persistent embryonic veins), and soft tissue or bone hypertrophy of the affected limb (typically a lower extremity). Caused by a somatic activating mutation in PIK3CA (part of the PIK3CA-related overgrowth spectrum). The hypertrophy is disproportionate, involving bone (increased length and girth) and soft tissues. Present at birth; progresses with growth. Complications: chronic venous insufficiency, cellulitis, deep vein thrombosis, pulmonary embolism (from anomalous veins), and life-threatening bleeding from pelvic or abdominal venous malformations. Diagnosis: clinical and MRI/MRV of the affected limb (shows abnormal venous anatomy and hypertrophy). Treatment: compression stockings for venous insufficiency; epiphyseal stapling or surgical debulking for limb overgrowth; sclerotherapy or surgical excision for symptomatic venous malformations. Anticoagulation for thromboembolic events. Differentiated from Parkes Weber syndrome by the absence of arteriovenous fistulas.

\medterm{Legg-Calve-Perthes Disease}-- Idiopathic avascular necrosis of the femoral head in children, typically ages 4-8 years. More common in boys (4:1) and Caucasian children. Presents with limping (antalgic gait), hip/knee pain, and limited hip range of motion (especially abduction and internal rotation). Etiology: temporary disruption of blood supply to the femoral head. Diagnosis: hip X-ray shows fragmentation, flattening, and collapse of the femoral head (staging: Waldenstrom). MRI for early diagnosis. Prognosis: better in younger children (<6), with preservation of femoral head sphericity. Treatment: observation for mild cases; containment therapy (Petrie cast, abduction brace, or femoral/pelvic osteotomy) for moderate to severe.

\medterm{Low Birth Weight}-- Birth weight <2500 g, often due to prematurity and/or intrauterine growth restriction; further classified: very low (<1500 g), extremely low (<1000 g), and small for gestational age. Risks: increased neonatal mortality, hypothermia, hypoglycemia, respiratory distress, infection, feeding difficulties, and long-term growth and neurodevelopmental problems. Causes: preterm birth, placental insufficiency, maternal hypertension/preeclampsia, smoking, malnutrition, multiple pregnancy, infection, and congenital anomalies. Management: specialized neonatal care (thermal, nutritional—breast milk/fortified feeds, infection prevention), monitoring for complications, catch-up growth, and long-term follow-up. Prevention: antenatal care and treatment of maternal conditions.

\medterm{Malrotation}-- See \hyperlink{Intestinal Malrotation}{Intestinal Malrotation}.

\medterm{Maple Syrup Urine Disease}-- An autosomal recessive inborn error of branched-chain amino acid metabolism caused by deficiency of the branched-chain alpha-ketoacid dehydrogenase complex (BCKDC) — mutations in BCKDHA, BCKDHB, DBT, or DLD. Results in accumulation of leucine, isoleucine, and valine (and their ketoacids) in the blood, giving a characteristic maple syrup (burnt sugar) odor to urine and cerumen. The classic neonatal form presents at 3-5 days of life with poor feeding, vomiting, lethargy, hypertonicity (opisthotonus), seizures, and cerebral edema progressing to coma and death if untreated. Diagnosis: plasma amino acid analysis (elevated leucine, isoleucine, valine, and alloisoleucine — pathognomonic) and BCKDHA/B gene testing. Newborn screening detects MSUD in most developed countries. Treatment: acute crises — rapid removal of branched-chain amino acids via hemodialysis or hemofiltration, IV glucose, insulin, and supplementation with isoleucine and valine to promote protein synthesis. Chronic management: strict dietary restriction of leucine, isoleucine, and valine (specialized medical formulas), thiamine supplementation (thiamine-responsive forms, 5\%), and metabolic monitoring. Liver transplantation can be curative.

\medterm{Marfanoid Habitus}-- A body habitus characterized by tall stature, long limbs, long fingers (arachnodactyly), and joint laxity. Seen in Marfan syndrome, but also in other connective tissue disorders and homocystinuria. Differentiated from Marfan syndrome by absence of aortic root dilation and lens subluxation. Requires cardiac evaluation to rule out aortic pathology.

\medterm{Measles}-- A highly contagious viral illness caused by Morbillivirus, transmitted by respiratory droplets. Prodrome of fever, coryza, cough, and conjunctivitis, followed by Koplik spots (white buccal mucosal lesions) and maculopapular rash spreading cephalocaudally. Complications: pneumonia, encephalitis, and subacute sclerosing panencephalitis. Prevention by MMR vaccination. Highly contagious with R0 of 12-18.

\medterm{Meconium Aspiration Syndrome}-- Respiratory distress in newborns born through meconium-stained amniotic fluid, caused by aspiration of meconium into the airways. Presents with respiratory distress, barrel-shaped chest, and coarse crackles. Meconium is thick, greenish-black fetal stool passed in utero. Risk factors include post-term pregnancy and fetal distress. Management: airway suctioning, mechanical ventilation, and treatment of complications including pneumothorax.

\medterm{Meconium-Stained Amniotic Fluid}-- Greenish or yellowish discoloration of amniotic fluid due to fetal meconium passage in utero. Associated with fetal distress, post-term pregnancy, and fetal growth restriction. Thick meconium increases risk of meconium aspiration syndrome. Management of non-vigorous newborns: tracheal suctioning before positive pressure ventilation. Vigorous newborns: routine care with close monitoring.

\medterm{Meningococcal Vaccine}-- Conjugate vaccine (MenACWY) and serogroup B vaccine (MenB) given at 11-12 years with booster at 16 years. Protects against Neisseria meningitidis, causing meningitis and septicemia. MenACWY booster at 16 years covers the high-risk period for serogroup C and Y disease. MenB given based on shared clinical decision-making or during outbreaks. One of the most devastating pediatric infections.

\medterm{Menkes Disease}-- An X-linked recessive disorder of copper metabolism caused by ATP7A mutations, resulting in impaired intestinal copper absorption and defective copper transport across the blood-brain barrier. Copper accumulates in intestinal and renal cells but is deficient in the liver and brain. Presents at 2-3 months with progressive neurologic deterioration, seizures, and characteristic hair abnormalities — stubby, white, brittle, and twisted (pili torti, pathognomonic on microscopy). Other features: hypotonia, failure to thrive, metaphyseal spurring (bladder-like metaphyses), wormian skull bones, and arterial tortuosity (cerebral and femoral arteries, prone to thrombosis and aneurysms). Diagnosis: low serum copper and ceruloplasmin (after 6-8 weeks of age), elevated copper in cultured fibroblasts, and ATP7A sequencing. Treatment: early parenteral copper-histidine supplementation (before symptoms) can improve neurologic outcomes. Established neurologic disease is irreversible. Prognosis: severe neurodegeneration and death by 3 years without treatment.

\medterm{Metaphyseal Corner Fractures}-- Bucket-handle or corner fractures at the metaphysis of long bones, characteristic of child abuse. Caused by twisting or pulling forces on extremities. Most common in the knee, ankle, and elbow. Seen on skeletal survey; may not be visible initially. Highly specific for non-accidental trauma in infants. Orthopedic consultation and full skeletal survey recommended.

\medterm{Micrognathia}-- A condition where the mandible is abnormally small, resulting in a retruded chin and decreased size of the oral cavity. May be isolated or syndromic (Pierre Robin sequence, Treacher Collins, Nager syndrome, Stickler syndrome). Presents at birth with potential airway obstruction (glossoptosis) and feeding difficulties. The severity depends on the degree of mandibular hypoplasia and tongue position. Diagnosis: clinical examination and cephalometric analysis. Management: prone positioning or lateral positioning for mild cases (gravity helps keep the tongue forward), nasopharyngeal airway for moderate obstruction, and mandibular distraction osteogenesis for severe obstruction. Most cases show significant mandibular catch-up growth during childhood.

\medterm{MMR Vaccine}-- Measles, Mumps, and Rubella live attenuated vaccine given at 12-15 months and 4-6 years. Two doses provide 97\% protection against measles and mumps, and 98\% against rubella. Contraindicated in pregnancy and immunocompromised individuals. Post-exposure prophylaxis within 72 hours. Measles remains a global health threat; outbreaks occur in under-vaccinated communities.

\medterm{Measles-Rubella (MR) Vaccine}-- A combined live attenuated vaccine containing measles and rubella antigens without the mumps component, widely used in countries where mumps is not a significant public health burden, including India's Universal Immunization Schedule. The WHO-prequalified MR vaccine (e.g., Serum Institute of India) is administered as two doses: the first at 9 months (alongside or shortly after BCG) and the second between 16-24 months. In India, the first dose is given at 9-12 months and the second at 16-24 months. Two doses provide greater than 95 percent protection against both measles and rubella. Measles remains a leading cause of vaccine-preventable childhood mortality globally, and rubella infection in pregnancy causes congenital rubella syndrome (deafness, cataracts, cardiac defects). The MR vaccine replaces the separate monovalent measles and rubella doses used in earlier schedules. Contraindicated in pregnancy and severely immunocompromised individuals. Post-exposure prophylaxis within 72 hours. See also \hyperlink{MMR Vaccine}{MMR Vaccine} for the mumps-containing formulation.

\medterm{Morgagni Hernia}-- A congenital diaphragmatic hernia occurring through the anterior retrosternal foramen of Morgagni (space of Larrey), resulting from failure of the sternal and costal portions of the diaphragm to fuse. The right side is affected in 90\% of cases (the hernia sac is present in 90\%). Abdominal contents (typically omentum, transverse colon, and liver) herniate into the anterior mediastinum. Often asymptomatic and incidentally discovered on chest X-ray (seen as a right cardiophrenic angle mass). Presents later in life than Bochdalek hernias (childhood to adulthood) with recurrent chest infections, dyspnea, or abdominal pain. Associated with obesity and trauma (increases intra-abdominal pressure). Diagnosis: chest X-ray shows a well-defined opacity in the right cardiophrenic angle; CT confirms the hernia, the contents, and the defect. Treatment: surgical repair (laparoscopic or open) with reduction of hernia contents and primary closure of the defect with or without mesh. Prognosis: excellent, with low recurrence rates.

\medterm{Moro Reflex}-- A primitive reflex present from birth to 3-6 months, elicited by suddenly startling the infant (loud noise, head drop). The infant extends both arms outward with fingers spread, then draws arms inward as if embracing. Absence at birth suggests central nervous system injury or brachial plexus palsy. Asymmetric response may indicate clavicle fracture or Erb palsy. Used as a marker of neurologic integrity.

\medterm{Multicystic Dysplastic Kidney}-- A congenital renal anomaly where one or both kidneys are replaced by multiple non-communicating cysts of varying sizes, with no functioning renal parenchyma. The most common cause of an abdominal mass in a newborn (unilateral). The ureter is typically atretic. The contralateral kidney may have VUR (20\%) or UPJ obstruction (10\%). Bilateral MCDK is incompatible with life (Potter syndrome). Most unilateral cases are asymptomatic and discovered on prenatal ultrasound. Diagnosis: postnatal ultrasound shows multiple non-communicating cysts of varying sizes, no identifiable renal sinus, and a reniform shape may be lost. A DMSA renal scan shows no function in the affected kidney. Treatment: observation is standard — most MCDKs involute (shrink) spontaneously over years. Blood pressure monitoring and urinalysis annually. Nephrectomy is rarely indicated (hypertension, recurrent UTI, or mass effect). Contralateral kidney evaluation (VCUG) for VUR and (US) for UPJ obstruction is essential.

\medterm{Mumps}-- Paramyxovirus infection causing bilateral parotitis, fever, and headache. Complications include orchitis (rarely causing infertility), meningitis, pancreatitis, and sensorineural hearing loss. Prevention by MMR vaccination. Outbreaks occur in unvaccinated populations. Diagnosis by serology or RT-PCR from saliva. Supportive treatment; no specific antiviral therapy.

\medterm{Myelomeningocele}-- The most severe form of spina bifida, where the spinal cord and meninges herniate through a vertebral defect into a sac on the back. Presents with a visible, fluid-filled sac at birth. Associated with hydrocephalus, Arnold-Chiari malformation, and lower extremity paralysis. Management: surgical closure within 24-48 hours, CSF drainage if hydrocephalus develops, and long-term rehabilitation for mobility and bowel/bladder function.

\medterm{Necrotizing Enterocolitis}-- An acquired necrotic inflammatory disease of the intestine, primarily affecting premature infants. Risk factors include prematurity, enteral feeding, and bacterial colonization. Presents with abdominal distension, bloody stools, bilious vomiting, and feeding intolerance. Pneumatosis intestinalis on X-ray is pathognomonic. Management: NPO, nasogastric decompression, IV antibiotics, and supportive care. Complications include perforation, peritonitis, and stricture formation.

\medterm{Neglect}-- A form of child maltreatment: the failure to provide for a child's basic needs—physical (food, clothing, shelter, medical care), educational, or emotional—causing harm or risk of harm. Features: failure to thrive, poor hygiene, untreated illness/injury, missed healthcare appointments, inadequate supervision, and developmental delay. Risk: parental mental illness, substance use, poverty, and domestic violence. Diagnosis: high index of suspicion, history, physical exam, growth assessment, and documentation. Management: safeguarding—report to child protection services, multidisciplinary assessment, support for the family (parenting programs, social services), and child safety planning; hospitalization for severe cases. Prevention: home visiting and family support.

\medterm{Neonatal Abstinence Syndrome}-- A withdrawal syndrome in neonates following in utero exposure to opioids (most common) or other substances (benzodiazepines, SSRIs, alcohol). Onset within 24-72 hours of birth (varies by half-life of substance). Presents with CNS hyperirritability (high-pitched crying, tremors, irritability, hypertonia, seizures), autonomic dysfunction (sweating, fever, tachypnea), and GI symptoms (poor feeding, vomiting, diarrhea). Diagnosis: Finnegan Neonatal Abstinence Scoring System quantifies severity. Treatment: non-pharmacologic (swaddling, low stimulation, skin-to-skin contact, breastfeeding if no contraindications) and pharmacologic for moderate-severe (morphine or methadone weaning protocol). Adjunctive: clonidine, phenobarbital. Prognosis: most recover fully with appropriate management.

\medterm{Neonatal Cephalohematoma}-- A subperiosteal hematoma of the skull, occurring in approximately 1-2\% of births, usually from forceps or vacuum delivery. Presents as a boggy, fluctuant swelling over the skull that does not cross suture lines (differentiating from caput succedaneum). Resolves spontaneously over weeks to months. Complications include hyperbilirubinemia, anemia, and calcification. No treatment required in most cases.

\medterm{Neonatal Herpes Simplex}-- Herpes simplex virus infection in the first 28 days of life, usually HSV-2 from maternal genital infection. Three forms: skin, eye, and mouth disease; CNS disease; and disseminated disease. Disseminated form mimics sepsis with high mortality. Diagnosis by viral PCR, culture, or Tzanck smear. Treatment: IV acyclovir for 14-21 days. Cesarean delivery recommended for active genital lesions near term.

\medterm{Neonatal Hypoglycemia}-- Blood glucose less than 40-45 mg/dL in term infants or less than 25 mg/dL in preterm infants. Risk factors include prematurity, small for gestational age, large for gestational age, and maternal diabetes. Presents with jitteriness, lethargy, poor feeding, and seizures. Treatment: early feeding, IV dextrose infusion, and monitoring. Persistent hypoglycemia may require glucagon or hydrocortisone.

\medterm{Neonatal Hypothermia}-- Core body temperature below 36.5 degrees Celsius in a newborn. Classified as mild (36-36.4), moderate (32-35.9), or severe (below 32). Occurs due to rapid heat loss through evaporation, convection, conduction, and radiation. Prevention: skin-to-skin contact, drying, warm environment, and delayed bathing. Treatment: warming methods appropriate to severity; monitor glucose and electrolytes.

\medterm{Neonatal Jaundice}-- See \hyperlink{Neonatal Jaundice (Jaundice of the Newborn)}{Neonatal Jaundice (Jaundice of the Newborn)}.

\medterm{Neonatal Jaundice (Jaundice of the Newborn)}-- Yellowish discolouration of the skin and sclerae in newborn infants due to elevated unconjugated bilirubin. Occurs in 60\% of term infants and 80\% of preterm infants. Physiological jaundice: appears after 24 hours of life, peaks at day 3-5 in term infants (day 5-7 in preterm), resolves by 2 weeks. Caused by: increased RBC mass, shorter RBC lifespan (70-90 days), immature hepatic conjugation (decreased UGT activity), increased enterohepatic circulation. Pathological jaundice: appears within 24 hours, bilirubin rising >85 micromol/L/day, or persisting >14 days (term) or >21 days (preterm). Causes: hemolytic (ABO incompatibility, Rh incompatibility, G6PD deficiency, hereditary spherocytosis), bruising/cephalhaematoma, polycythaemia, sepsis, biliary atresia, galactosaemia, Crigler-Najjar syndrome, Gilbert syndrome. Complications: acute bilirubin encephalopathy (lethargy, poor feeding, hypotonia, high-pitched cry, opisthotonos, seizures) and kernicterus (chronic, irreversible choreoathetoid cerebral palsy, upward gaze palsy, sensorineural hearing loss, enamel dysplasia). Management: phototherapy (most effective, blue light 460-490 nm, converts bilirubin to water-soluble isomers excreted without conjugation), exchange transfusion (for severe hyperbilirubinaemia failing phototherapy), and treatment of underlying cause. NICE and AAP guidelines provide phototherapy and exchange transfusion thresholds based on gestational age, birth weight, and risk factors.

\medterm{Neonatal Sepsis}-- Systemic bacterial infection in the first 28 days of life, classified as early-onset (within 72 hours, usually from vertical transmission) or late-onset (after 72 hours, often nosocomial). Common organisms: Group B Streptococcus, E. coli, Listeria. Presents with non-specific symptoms: poor feeding, lethargy, temperature instability, tachycardia, and respiratory distress. Diagnosis by blood culture; treatment with empiric IV antibiotics (ampicillin and gentamicin).

\medterm{Neonatologist}-- A pediatrician who specializes in the medical care of newborn infants, particularly premature and critically ill neonates. Neonatologists manage conditions including respiratory distress syndrome, bronchopulmonary dysplasia, intraventricular hemorrhage, necrotizing enterocolitis, neonatal sepsis, hypoxic-ischemic encephalopathy (therapeutic hypothermia), congenital anomalies, and extreme prematurity. They work in neonatal intensive care units (NICU), lead resuscitation teams at high-risk deliveries, and coordinate care with pediatric surgeons, cardiologists, and geneticists. Training: pediatric residency + 3-year neonatal-perinatal medicine fellowship.

\medterm{Neural Tube Defects}-- Congenital malformations resulting from failure of the neural tube to close during the third to fourth week of gestation. Include anencephaly (incomplete cranial closure), encephalocele (herniation through skull defect), and spina bifida (incomplete spinal closure). Spina bifida ranges from occulta (mild, asymptomatic) to myelomeningocele (severe, with neural tissue protrusion). Prevention with folic acid supplementation before and during pregnancy.

\medterm{Neurofibromatosis Type 2}-- An autosomal dominant disorder caused by NF2 gene mutation on chromosome 22. Characterized by bilateral vestibular schwannomas (acoustic neuromas), meningiomas, and ependymomas. Features: hearing loss, tinnitus, and balance problems. Diagnosis: MRI showing bilateral vestibular schwannomas. Management: hearing preservation surgery, radiation therapy, and surveillance imaging.

\medterm{Newborn Circumcision}-- Surgical removal of the foreskin from the penis, typically performed within the first few days of life. Medical benefits include reduced risk of urinary tract infections, sexually transmitted infections, and penile cancer. Complications include bleeding, infection, and meatal stenosis. Techniques include Plastibell, Gomco clamp, and Mogen clamp. Performed with parental informed consent. Not universally recommended; benefits versus risks should be discussed.

\medterm{Niemann-Pick Disease}-- A group of autosomal recessive lysosomal storage disorders caused by sphingomyelinase deficiency (ASM, types A and B, SMPD1 mutations) or NPC1/NPC2 cholesterol transport defects (type C). Type A (infantile neuropathic): severe, presents at 3-6 months with hepatosplenomegaly, cherry-red macula (50\%), developmental regression, neurodegeneration, and death by 2-3 years. Type B (chronic visceral): later onset, hepatosplenomegaly, pulmonary infiltrates (foam cells in alveoli), thrombocytopenia, and dyslipidemia — no or minimal neurologic involvement; survival into adulthood. Type C: variable onset (infantile to adult) with neurologic features: vertical supranuclear gaze palsy (pathognomonic), ataxia, dystonia, seizures, dementia, and cataplexy. Diagnosis: types A/B — acid sphingomyelinase enzyme assay in leukocytes. Type C — filipin staining for cholesterol accumulation in fibroblasts and NPC1/NPC2 sequencing. Treatment: types A/B — olipudase alfa (enzyme replacement therapy, FDA-approved for type A/B non-neurologic). Type C — miglustat (substrate reduction therapy). Supportive care for neurological symptoms.

\medterm{Non-Accidental Trauma}-- Injuries inflicted on a child, suspected when history and examination findings are inconsistent with the reported mechanism. Red flags: multiple fractures at different stages of healing, posterior rib fractures, metaphyseal corner fractures, and retinal hemorrhages. Requires comprehensive evaluation including skeletal survey, CT head, ophthalmologic exam, and laboratory workup. Social work and child protective services involvement mandatory.

\medterm{Noonan Syndrome}-- An autosomal dominant multisystem disorder caused by mutations in the RAS-MAPK pathway genes (PTPN11 50\%, SOS1 10-15\%, RAF1 5-10\%, RIT1, KRAS). Features: characteristic facies (hypertelorism, downslanting palpebral fissures, low-set posteriorly rotated ears, webbed neck, ptosis), short stature, congenital heart disease (pulmonary stenosis 50-60\%, hypertrophic cardiomyopathy 20-30\%), pectus deformities, and cryptorchidism. Coagulation defects (factor XI deficiency, von Willebrand disease) are common. Intelligence is usually normal, though learning disabilities and speech delays are common. Diagnosis: clinical criteria (van der Burgt) and genetic testing. Management: growth hormone therapy for short stature, cardiology follow-up (pulmonary stenosis balloon valvuloplasty, HCM monitoring), and developmental support.

\medterm{Omphalitis}-- Infection of the umbilical stump in neonates, typically occurring within the first 1--2 weeks of life. Omphalitis is a potentially serious condition requiring prompt treatment due to the risk of spread to the liver (portal vein thrombosis, portal hypertension) and peritoneum via the umbilical vessels. Risk factors: home delivery with unsterile cord care, low birth weight, prolonged rupture of membranes (>18 hours), umbilical catheterisation, and immunocompromised state. Presentation: periumbilical erythema, edema, tenderness, purulent or foul-smelling discharge from the umbilicus. Signs of systemic infection (fever, lethargy, poor feeding, tachycardia, hypotension) indicate severe disease. Pathogens: Staphylococcus aureus, Group A Streptococcus, Escherichia coli, Klebsiella, and anaerobes (Bacteroides fragilis). Diagnosis is clinical; wound culture, blood culture, and ultrasound (assess for abscess, thrombophlebitis, necrotising fasciitis) in severe cases. Treatment: IV antibiotics covering Gram-positive, Gram-negative, and anaerobic organisms (flucloxacillin plus gentamicin plus metronidazole). Surgical debridement for necrotising fasciitis or abscess drainage. Complications: necrotising fasciitis (abdominal wall cellulitis with crepitus, high mortality), portal vein thrombosis, peritonitis, and neonatal sepsis. Prevention: chlorhexidine cord care (reduces omphalitis in high-mortality settings per WHO recommendation).

\medterm{Omphalocele}-- A congenital abdominal wall defect where abdominal organs herniate into the base of the umbilical cord, covered by a peritoneal-amniotic membrane. Associated with chromosomal abnormalities (trisomies 13, 18) and other anomalies. Differentiated from gastroschisis by membrane coverage and midline location. Small omphaloceles may be repaired primarily; large defects require staged closure with silastic silo and gradual reduction.

\medterm{Oral Polio Vaccine (OPV-O)}-- Sabin vaccine, a live attenuated preparation of poliovirus administered orally, historically the primary tool for global polio eradication. The bivalent formulation (bOPV) contains attenuated types 1 and 3 only; the type 2 component was withdrawn from routine use in April 2016 following the certified eradication of wild poliovirus type 2. OPV replicates in the intestinal mucosa, inducing both humoral (serum IgG) and robust mucosal (secretory IgA) immunity, which interrupts person-to-person fecal-oral transmission---a critical advantage over inactivated polio vaccine (IPV) for eradication campaigns. Advantages include low cost, ease of administration (no needles), cold-chain flexibility, and secondary immunisation of close contacts through shedding of attenuated virus. The World Health Organization and Global Polio Eradication Initiative continue to use bOPV in mass vaccination campaigns in endemic and at-risk regions. Risks include vaccine-associated paralytic polio (VAPP, approximately 1 case per 2.7 million first doses, caused by reversion of attenuated virus to neurovirulent form) and circulating vaccine-derived poliovirus (cVDPV), which emerges in under-immunised populations where the attenuated virus circulates and reverts. OPV is contraindicated in immunocompromised individuals. The United States exclusively uses IPV since 2000. In India's Universal Immunization Schedule, OPV-0 is given at birth, followed by OPV 1-3 at 6, 10, and 14 weeks, and a booster at 16-24 months, alongside fractional IPV for enhanced protection. See \hyperlink{India Immunization Schedule}{India Immunization Schedule}, \hyperlink{Polio}{Polio}, and \hyperlink{Poliomyelitis}{Poliomyelitis} in infectious disease.

\medterm{Osteogenesis Imperfecta}-- A connective tissue disorder caused by collagen type I mutations, ranging from mild to lethal. Features: fractures with minimal trauma, blue sclerae, dental abnormalities, and hearing loss. Severe forms (type II) cause multiple fractures in utero and perinatal death. Diagnosis: clinical features, genetic testing, and bone biopsy. Management: bisphosphonates, physical therapy, and surgical correction of deformities.

\medterm{Otitis Media}-- Middle ear infection, the most common childhood illness after the common cold. Peak incidence between 6-15 months. Risk factors include secondhand smoke, daycare attendance, formula feeding, and anatomy. Presents with ear pain, fever, irritability, and hearing loss. Diagnosis by otoscopic exam showing erythema, bulging tympanic membrane, or effusion. Treatment: antibiotics (amoxicillin first-line) for acute otitis media. See also \hyperlink{Acute Otitis Media}{Acute Otitis Media}.

\medterm{Palmar Grasp Reflex}-- A primitive reflex present from birth to 5-6 months, elicited by placing a finger in the infant's palm. The infant grips tightly, strong enough to support the infant's weight briefly. Bilateral symmetry assessed; asymmetry may indicate birth injury. Disappears as voluntary grasping develops around 3-4 months. Used as a neurologic marker in the newborn examination.

\medterm{Parvovirus B19}-- Human parvovirus B19: a small DNA virus causing erythema infectiosum ('fifth disease') in children, transmitted by respiratory droplets. Features: 'slapped cheek' rash (bright red cheeks) followed by a lace-like (reticular) rash on the trunk/limbs, mild fever; in adults: arthropathy; in chronic hemolytic states (sickle cell): transient aplastic crisis (severe anemia); in pregnancy: fetal hydrops (anemia) and fetal loss; in immunocompromised: chronic anemia. Diagnosis: clinical; serology (IgM/IgG), PCR (aplastic crisis/fetus). Management: supportive—transfusion for aplastic crisis, intrauterine transfusion for severe fetal anemia; no vaccine; respiratory precautions; reassure about benign course in immunocompetent children.

\medterm{Patau Syndrome}-- Trisomy 13, affecting 1 in 10,000-20,000 live births. Features: holoprosencephaly (most characteristic, incomplete forebrain division), microcephaly, microphthalmia or anophthalmia, cleft lip and palate, polydactyly (postaxial), clenched fists with overlapping fingers, and congenital heart defects (VSD, PDA, dextrocardia, 80\%). Other findings: cutis aplasia (scalp defect, particularly the occipital scalp), omphalocele, and renal anomalies (polycystic kidneys). Diagnosis: prenatal screening (abnormal first-trimester ultrasound with increased nuchal translucency), NIPT, and karyotype/CMA. Management: comfort care (median survival 7-10 days; <10\% survive to 1 year). Survivors have profound intellectual disability and severe growth restriction.

\medterm{Patent Ductus Arteriosus}-- PDA: failure of the ductus arteriosus (fetal connection between pulmonary artery and aorta) to close after birth, causing a left-to-right shunt. Common in preterm infants; term—often small. Features: continuous 'machinery' murmur, bounding pulses, wide pulse pressure; large shunt—heart failure, failure to thrive, pulmonary hypertension (Eisenmenger if long-standing). Diagnosis: echocardiography. Management: preterm—indomethacin/ibuprofen (cyclooxygenase inhibitors) or paracetamol, and surgical ligation if refractory; term/older—transcatheter closure (coil/device) or surgery; small asymptomatic PDAs—monitoring/elective closure; prevention of endocarditis historically. Closure prevents complications.

\medterm{Pectus Carinatum}-- A congenital chest wall deformity characterized by anterior protrusion of the sternum (pigeon chest), caused by excessive growth of the costal cartilages pushing the sternum forward. Less common than pectus excavatum. Associated with congenital heart disease, Marfan syndrome, Noonan syndrome, and homocystinuria. Presents in childhood and progresses during the adolescent growth spurt. Symptoms are predominantly cosmetic, but some children experience chest pain, exercise intolerance, or psychosocial distress. Diagnosis: clinical examination and CT or X-ray to evaluate symmetry and severity. Treatment: orthotic bracing (dynamic chest compressor, effective for flexible deformities in younger children) is first-line, worn 12-23 hours daily for 6-12 months. Surgical correction (Ravitch-type sternal osteotomy with costal cartilage resection) for severe or rigid deformities or bracing failure. Outcomes are excellent with both methods.

\medterm{Pectus Excavatum}-- The most common congenital chest wall deformity, characterized by depression of the sternum (funnel chest), caused by excessive growth of the costal cartilages that push the sternum inward. Typically sporadic, but associated with Marfan syndrome, Ehlers-Danlos syndrome, Noonan syndrome, and Poland syndrome. Presents in early childhood and progresses during the adolescent growth spurt. Symptoms are usually cosmetic, but severe cases can cause cardiopulmonary impairment (reduced stroke volume, dyspnea on exertion, exercise intolerance). The Haller index (CT ratio of transverse chest diameter to anterior-posterior diameter, >3.25 is severe) quantifies severity. Diagnosis: clinical (the depression is obvious) and CT chest for Haller index. Treatment: observation for mild cases. Surgical correction: Nuss procedure (minimally invasive, pectus bar placed retrosternally under thoracoscopic guidance, removed after 2-4 years) for moderate to severe cases. Modified Ravitch procedure (open resection of costal cartilages and sternal osteotomy) for older patients or failed Nuss. Outcomes: excellent cosmetic and functional results.

\medterm{Pediatric Diabetic Ketoacidosis}-- A life-threatening complication of type 1 diabetes characterized by hyperglycemia, metabolic acidosis, and ketosis. Presents with polyuria, polydipsia, vomiting, abdominal pain, Kussmaul respirations, and fruity breath odor. Diagnosis: blood glucose greater than 250 mg/dL, pH less than 7.3, bicarbonate less than 15 mEq/L, and positive ketones. Treatment: IV fluids, insulin infusion, and electrolyte replacement with careful potassium monitoring.

\medterm{Pediatric Health}-- Pediatric health encompasses the physical, cognitive, emotional, and social development of children from infancy through adolescence, with emphasis on preventive care, developmental surveillance, and age-appropriate management of acute and chronic conditions. Well-child care follows a schedule: immunizations (per CDC/ACIP: DTaP, IPV, MMR, varicella, HPV, meningococcal, influenza, COVID-19), developmental screening (Ages and Stages Questionnaire, M-CHAT for autism at 18 and 24 months), growth monitoring (WHO/CDC growth charts), vision/hearing screening, and anticipatory guidance (injury prevention, nutrition, sleep, media use, substance use prevention). Common pediatric conditions: otitis media, asthma, atopic dermatitis, allergic rhinitis, ADHD (prevalence ~9\%), autism spectrum disorder (prevalence ~2.8\%), anxiety/depression, obesity (affecting ~20\% of US children), infectious diseases (RSV, croup, gastroenteritis), and congenital anomalies. Adolescent health addresses puberty, sexual health (STI prevention, contraception), mental health (leading cause of disability in teens), substance use, and transition to adult care. Family-centered care and trauma-informed approaches improve outcomes.

\medterm{Pediatric Meningitis}-- Inflammation of the meninges in children, bacterial forms being the most dangerous. Presents with fever, headache, neck stiffness, photophobia, altered mental status, and petechial rash (meningococcal). Infants may present with irritability, poor feeding, bulging fontanelle, and high-pitched cry. Diagnosis by lumbar puncture showing elevated WBC, elevated protein, low glucose, and positive Gram stain/culture. Treatment: empiric ceftriaxone and vancomycin.

\medterm{Pediatrician}-- A medical doctor who specializes in the health and medical care of children from birth through adolescence (typically up to age 16-18). Pediatricians focus on: growth and development, preventive care (immunizations, developmental screening, anticipatory guidance), acute illness (infections, injuries), chronic disease management (asthma, diabetes, epilepsy, allergies), and behavioral/mental health. Subspecialties: neonatology, pediatric cardiology, pediatric oncology, pediatric gastroenterology, pediatric emergency medicine, pediatric intensive care, pediatric endocrinology, pediatric neurology, and developmental pediatrics. Training: medical school + 2 years foundation training (UK) or 3 years pediatrics residency (US) + subspecialty fellowship.

\medterm{Pediatrics}-- The branch of medicine concerned with the health and medical care of infants, children, and adolescents from birth to age 18 (or up to age 16 in the UK). Pediatric care encompasses: well-child visits (growth monitoring, developmental screening, immunisations, anticipatory guidance), management of acute illness (infections, injuries), management of chronic conditions (asthma, diabetes, epilepsy, cystic fibrosis, congenital heart disease), and behavioural/mental health (ADHD, autism, anxiety, depression). Subspecialties: neonatology, pediatric cardiology, pediatric oncology, pediatric gastroenterology, pediatric pulmonology, pediatric neurology, pediatric nephrology, pediatric endocrinology, pediatric rheumatology, pediatric infectious disease, pediatric emergency medicine, pediatric intensive care, developmental pediatrics, and pediatric surgery. Preventive care: the childhood immunisation schedule (UK: 2, 3, 4 months; 12--13 months; 3 years 4 months; 12--14 years—boosters for tetanus, diphtheria, polio, HPV, MenACWY). The UK has a comprehensive child health surveillance programme (health visitor reviews, school entry checks).

\medterm{Periventricular Leukomalacia}-- Ischemic white matter injury in the premature infant brain, a major cause of cerebral palsy. Risk factors include prematurity, hypotension, and inflammation. Diagnosis by cranial ultrasound showing periventricular echodensities or cystic changes. Prevention: avoiding hypotension, maintaining normothermia, and managing infection. Long-term sequelae include spastic diplegia, cognitive impairment, and visual disturbances.

\medterm{Pertussis}-- A highly contagious bacterial infection caused by Bordetella pertussis, producing characteristic paroxysmal cough with whooping inspiratory sound. Three stages: catarrhal (1-2 weeks of URI symptoms), paroxysmal (weeks to months of violent coughing spells), and convalescent. Most dangerous in infants under 6 months. Treatment: macrolide antibiotics (azithromycin). Prevention: DTaP vaccination.

\medterm{Pentavalent Vaccine}-- A WHO-prequalified five-in-one combination vaccine containing Diphtheria toxoid, Tetanus toxoid, whole-cell Pertussis antigen, Hepatitis B (HBsAg), and Haemophilus influenzae type b (Hib) conjugate, administered as a single intramuscular injection. The pentavalent formulation replaces separate DTaP, HepB, and Hib injections, reducing the number of injections and improving coverage in resource-limited settings. India's Universal Immunization Schedule uses pentavalent vaccine at 6, 10, and 14 weeks of age (three primary doses). The whole-cell pertussis component (as opposed to the acellular pertussis in US DTaP) provides broader protection but may cause more local reactions and fever. Common adverse effects include injection site swelling, redness, and pain, and transient fever (managed with paracetamol). The pentavalent vaccine is a cornerstone of the Expanded Programme on Immunization (EPI) in over 100 countries. Alternative formulations include DPT-HepB (without Hib) where Hib conjugate is supplied separately. DPT booster doses are given at 16-24 months and 5-6 years as part of India's Universal Immunization Schedule. See \hyperlink{India Immunization Schedule}{India Immunization Schedule}.

\medterm{Pfeiffer Syndrome}-- An autosomal dominant craniosynostosis syndrome caused by FGFR1 or FGFR2 mutations. Features: craniosynostosis (bilateral coronal), midface hypoplasia, proptosis, hypertelorism, broad thumbs and great toes (medial deviation), and variable syndactyly. Three clinical types: Type 1 (classic, milder, normal intelligence, good prognosis with surgery), Type 2 (cloverleaf skull, severe proptosis, elbow ankylosis, poor prognosis), and Type 3 (similar to Type 2 without cloverleaf skull, severe neurological impairment). Types 2 and 3 have significant mortality. Diagnosis: clinical and genetic testing. Management: staged cranial vault and midface surgery for Type 1; palliative management for Types 2 and 3.

\medterm{Pierre Robin Sequence}-- A congenital anomaly triad: micrognathia (small mandible), glossoptosis (posteriorly displaced tongue causing airway obstruction), and U-shaped cleft palate. The sequence arises from early mandibular hypoplasia, which prevents the tongue from descending, thereby impairing palatal shelf fusion. Presents at birth with respiratory distress (stridor, retractions, cyanosis, worse supine) and feeding difficulties. Diagnosis: clinical examination and polysomnography to assess obstructive sleep apnea severity. Management: prone positioning (first-line, 70\% improve), nasopharyngeal airway, and surgical intervention for severe cases (tongue-lip adhesion, mandibular distraction osteogenesis, or tracheostomy). Cleft palate repair is typically performed at 9-18 months.

\medterm{Pneumococcal Conjugate Vaccine}-- PCV13 and PCV15 vaccines given at 2, 4, 6, and 12-15 months. Protects against Streptococcus pneumoniae serotypes causing invasive disease, pneumonia, and otitis media. PCV20 now recommended for older children and adults. Has dramatically reduced rates of invasive pneumococcal disease in children since universal vaccination.

\medterm{Poland Syndrome}-- A congenital anomaly characterized by unilateral absence or hypoplasia of the pectoralis major muscle (specifically the sternocostal head) and ipsilateral hand anomalies (syndactyly, most commonly; also brachydactyly, ectrodactyly, or radial club hand). The right side is affected in 60-75\% of cases. Males are affected 2-3 times more frequently. Associated findings: rib defects (aplasia or hypoplasia of the second to fifth ribs), chest wall depression, breast and nipple hypoplasia or absence (in females), axillary hair loss, and vertebral anomalies. The cause is thought to be an interruption of the subclavian artery blood supply during the sixth week of gestation. Diagnosis: clinical with ultrasound or CT for chest wall evaluation. Treatment: hand surgery for syndactyly release; breast reconstruction with implants or autologous tissue for females; chest wall reconstruction for severe rib defects.

\medterm{Post-Term Pregnancy}-- Pregnancy continuing beyond 42 weeks gestation (294 days from LMP). Risk factors include prior post-term pregnancy, nulliparity, and genetic factors. Complications include macrosomia, oligohydramnios, meconium aspiration, placental insufficiency, and stillbirth. Management: antenatal surveillance with NST and biophysical profile starting at 41 weeks. Induction of labor recommended by 41-42 weeks.

\medterm{Posterior Urethral Valves}-- The most common cause of lower urinary tract obstruction in male infants, caused by congenital mucosal folds (valves) in the posterior urethra (prostatic urethra) that obstruct urine outflow. Presents prenatally with oligohydramnios, bilateral hydronephrosis, and a distended, thick-walled bladder, and postnatally with poor urinary stream, urinary retention, and abdominal distension. Severe cases present with the prune belly sequence (deficient abdominal wall musculature from massive bladder distension) and Potter sequence (pulmonary hypoplasia from oligohydramnios). Diagnosis: prenatal ultrasound, postnatal voiding cystourethrogram (VCUG) shows a dilated posterior urethra with a sharp cutoff. Treatment: immediate bladder decompression (urethral catheter), endoscopic valve ablation (transurethral incision of the valves, first-line). Temporary vesicostomy for small or premature infants. Prognosis: depends on the degree of renal impairment at presentation. Even with successful valve ablation, 25-30\% develop end-stage renal disease in childhood or adolescence due to preexisting renal dysplasia.

\medterm{Potter Syndrome}-- A sequence of anomalies resulting from severe oligohydramnios (low amniotic fluid volume), most commonly caused by bilateral renal agenesis, but also from autosomal recessive polycystic kidney disease, posterior urethral valves, or chronic amniotic fluid leak. Features: pulmonary hypoplasia (the most critical, incompatible with life), Potter facies (low-set, flattened ears, recessed chin, prominent epicanthal folds, flattened nose), limb deformities (clubfoot, hip dislocation, flexion contractures), and amnion nodosum (nodules on the amnion). Presents at birth with severe respiratory distress, inability to oxygenate, and typical facies. Diagnosis: prenatal ultrasound shows anhydramnios/oligohydramnios and absent or abnormal kidneys. Chest X-ray shows small lung volumes. Prognosis: uniformly fatal in the newborn period due to pulmonary hypoplasia. Management: perinatal palliative care.

\medterm{Prader-Willi Syndrome}-- A genetic disorder caused by loss of function of genes on chromosome 15 (paternal deletion or maternal uniparental disomy). Features: neonatal hypotonia and failure to thrive, followed by hyperphagia and obesity in childhood. Intellectual disability, short stature, hypogonadism, and behavioral problems. Management: controlled diet, growth hormone therapy, and behavioral interventions.

\medterm{Preterm Infant}-- An infant born before 37 completed weeks of gestation; classified by gestational age (late preterm 34-36, moderate 32-33, very preterm <32, extremely <28 weeks) and by birth weight. Causes: preterm labor, PPROM, preeclampsia, multiple pregnancy, infection, and iatrogenic. Complications: respiratory distress syndrome (surfactant deficiency), apnea, intraventricular hemorrhage, necrotizing enterocolitis, infection/sepsis, hypoglycemia, jaundice, thermoregulation difficulty, patent ductus arteriosus, and retinopathy of prematurity; long-term: neurodevelopmental and growth concerns. Management: neonatal intensive care (surfactant, respiratory support, thermal care, nutrition—breast milk, kangaroo care), prevention (progesterone, cervical cerclage), antenatal corticosteroids and magnesium for lung maturity/neuroprotection, and developmental follow-up.

\medterm{Prune Belly Syndrome}-- A congenital syndrome comprising the triad: deficient abdominal wall musculature (wrinkled, prune-like appearance), bilateral cryptorchidism (undescended testes), and urinary tract anomalies (dilated ureters, dilated and trabeculated bladder, and posterior urethral valves or urethral atresia). Almost exclusively affects boys. Other associations: pulmonary hypoplasia from oligohydramnios, orthopedic anomalies (clubfoot, hip dislocation), and cardiac anomalies. The severity ranges from the lethal form (severe oligohydramnios, pulmonary hypoplasia, and renal dysplasia) to the mild form (preserved renal function and mild abdominal wall laxity). Diagnosis: prenatal ultrasound shows a distended bladder with thickened walls and oligohydramnios. Postnatal diagnosis is clinical. Treatment: depends on severity — mild cases: observation and orchiopexy. Moderate-severe: urinary tract reconstruction (reduction cystoplasty, ureteral reimplantation, and abdominoplasty). Renal transplantation for end-stage renal disease. All patients require lifelong urologic and nephrologic surveillance.

\medterm{Reflux}-- Gastroesophageal reflux in infants: effortless regurgitation of gastric contents, very common (physiologic 'spitting up') and usually benign, resolving by 12 months with maturation of the lower esophageal sphincter. Pathologic GERD: poor weight gain, irritability, feeding refusal, arching (Sandifer syndrome), apnea/bradycardia, recurrent pneumonia (aspiration), and esophagitis. Diagnosis: clinical; endoscopy/PH-impedance for complications. Management: reassurance and feeding modifications (thickened feeds, frequent small feeds, upright positioning after feeds) for physiologic reflux; GERD—thickened feeds, alginate, and PPI (omeprazole) for significant disease; alarm signs (hematemesis, bilious vomiting, failure to thrive) require evaluation.

\medterm{Renal Agenesis}-- Congenital absence of one (unilateral, 1 in 1,000) or both (bilateral, 1 in 4,000) kidneys. Unilateral renal agenesis: usually asymptomatic, discovered incidentally — the solitary kidney undergoes compensatory hypertrophy. The contralateral kidney may have VUR, UPJ obstruction, or ectopia. Associated with VACTERL, Mayer-Rokitansky-Kuster-Hauser syndrome (Müllerian agenesis), and other GU anomalies. Bilateral renal agenesis (Potter syndrome): incompatible with life — causes severe oligohydramnios, leading to Potter sequence: pulmonary hypoplasia (the proximal cause of death), limb deformities (clubfoot, contractures), characteristic Potter facies (low-set ears, flat nose, epicanthal folds), and amnion nodosum. Diagnosis: prenatal ultrasound (absent renal tissue, empty renal fossa, severe oligohydramnios). Bilateral renal agenesis is uniformly fatal in the newborn period from pulmonary hypoplasia. Unilateral: annual blood pressure monitoring, urinalysis, and renal function surveillance.

\medterm{Respiratory Distress Syndrome}-- A condition primarily affecting premature infants caused by pulmonary surfactant deficiency, leading to alveolar collapse and respiratory failure. Presents with tachypnea, grunting, nasal flaring, retractions, and cyanosis. Chest X-ray shows diffuse ground-glass opacity. Treatment: surfactant replacement therapy, mechanical ventilation, and continuous positive airway pressure. Prevention with antenatal corticosteroids when preterm delivery anticipated.

\medterm{Respiratory Syncytial Virus}-- A paramyxovirus and the most common cause of bronchiolitis and pneumonia in infants and young children. Seasonal outbreaks in fall, winter, and spring. Transmission by respiratory droplets and fomites. Mild in older children and adults; can be severe or fatal in infants, premature infants, and those with chronic lung or heart disease. Prevention with palivizumab (Synagis) or nirsevimab (Beyfortus) in high-risk infants during RSV season. Nirsevimab is now the preferred RSV prophylaxis for all infants in many guidelines (CDC/ACIP), providing passive protection with a single dose.

\medterm{Retinal Hemorrhages}-- Bleeding within the retina, a hallmark finding in abusive head trauma (shaken baby syndrome). Found on dilated fundoscopic examination, often bilateral and multi-layered. Also seen in severe accidental trauma, bleeding disorders, and increased intracranial pressure. In the context of suspected abuse, strongly suggestive of non-accidental injury. Ophthalmologic consultation essential for documentation.

\medterm{Retinopathy of Prematurity}-- Abnormal vascularization of the retina in premature infants, a leading cause of childhood blindness. Affects infants less than 32 weeks gestation and less than 1500 grams. Stages range from mild (avascular retina) to severe (retinal detachment). Screening began at 31-33 weeks corrected gestational age or 4 weeks chronological age. Treatment: laser photocoagulation or anti-VEGF injection for type 1 ROP.

\medterm{Rett Syndrome}-- An X-linked dominant neurodevelopmental disorder affecting almost exclusively females, caused by mutations in the MECP2 gene. Normal development for the first 6-18 months, followed by regression with loss of acquired language and motor skills, stereotypical hand-wringing movements, deceleration of head growth (acquired microcephaly), seizures, and gait ataxia. Autonomic dysfunction causes breathing irregularities (hyperventilation, breath-holding), cold extremities, and bruxism. Diagnosis: clinical criteria (Neul 2010) and MECP2 genetic testing. Treatment: symptom management — antiepileptic drugs for seizures, physical/occupational therapy, nutritional support, and scoliosis monitoring. Trofinetide is FDA-approved for Rett syndrome in adults and children >2 years.

\medterm{Rheumatic Fever}-- A nonsuppurative inflammatory complication of Group A Streptococcal pharyngitis, occurring 2-4 weeks after infection. Diagnosed by Jones criteria: major manifestations (carditis, polyarthritis, chorea, erythema marginatum, subcutaneous nodules) plus evidence of prior streptococcal infection. Major criterion plus two minor criteria. Treatment: anti-inflammatory agents and antibiotics. Prevention: prompt treatment of strep throat.

\medterm{Rooting Reflex}-- A primitive reflex present from birth to 3-4 months, triggered by stroking the infant's cheek or corner of the mouth. The infant turns toward the stimulus and opens the mouth, searching for the nipple. Facilitates breastfeeding. Absence may indicate neurologic impairment. Disappears as voluntary feeding develops. One of several primitive reflexes assessed in the newborn period to evaluate neurologic function.

\medterm{Roseola}-- See \hyperlink{Roseola Infantum}{Roseola Infantum}.

\medterm{Roseola Infantum}-- A viral illness caused by Human Herpesvirus 6 (HHV-6), affecting children 6 months to 2 years. Characterized by high fever (39-40.5 degrees Celsius) for 3-5 days followed by defervescence and appearance of a maculopapular rash starting on the trunk and spreading outward. Rash appears as fever resolves. Diagnosis is clinical. Self-limited; treatment is supportive.

\medterm{Rotavirus Vaccine}-- Oral live attenuated vaccine given at 2, 4, and 6 months (RV1) or at 2, 4, and 8 months (RV5). Protects against rotavirus gastroenteritis, a leading cause of severe diarrhea in infants. Reduced rotavirus hospitalizations by over 80\% since introduction. Mild side effects include fever and diarrhea. Contraindicated in infants with severe combined immunodeficiency.

\medterm{RSV}-- Respiratory syncytial virus: the most common cause of bronchiolitis and viral pneumonia in infants and young children, transmitted by droplets/fomites; seasonal (winter). Features: coryza, cough, wheeze, tachypnea, and respiratory distress; apnea in young infants; risk factors for severe disease: prematurity, congenital heart disease, chronic lung disease, and immunodeficiency. Diagnosis: clinical; nasopharyngeal PCR/antigen. Management: supportive—oxygen, hydration/nasal suctioning, and respiratory support; no routine antiviral; palivizumab (monoclonal antibody) prophylaxis for high-risk infants; new vaccines/monoclonals (nirsevimab) emerging; prevention: hand hygiene, and avoid sick contacts.

\medterm{Rubella}-- German measles caused by Rubivirus. Mild illness with fever, rash, and lymphadenopathy in children. Dangerous in pregnancy due to congenital rubella syndrome (cataracts, heart defects, deafness). Rash similar to measles but milder, with postauricular and suboccipital lymphadenopathy. Prevention by MMR vaccination. Eradicated in the Americas since 2015.

\medterm{Rubinstein-Taybi Syndrome}-- A genetic disorder caused by mutations in CREBBP (50-60\%) or EP300 (3-8\%), encoding histone acetyltransferases involved in transcriptional regulation. Features: characteristic facies — downslanting palpebral fissures, broad nasal bridge, beaked nose, high-arched palate, and low-set ears; broad thumbs and halluces (pathognomonic, 100\%); short stature; microcephaly; and moderate to severe intellectual disability. Other findings: congenital heart disease (PDA, ASD, VSD, 30-35\%), cryptorchidism, renal anomalies, and increased risk of tumors (benign: meningioma, pilomatrixoma; malignant: leukemia). Diagnosis: clinical and genetic testing. Management: multidisciplinary — cardiac evaluation, developmental therapies, physical therapy for joint laxity, and tumor surveillance. Anesthesia risk: difficult airway from micrognathia and potential malignant hyperthermia.

\medterm{Scarlatina}-- A toxin-mediated illness caused by Group A Streptococcus producing erythrogenic toxin. See \hyperlink{Scarlet Fever}{Scarlet Fever}.

\medterm{Scarlet Fever}-- A toxin-mediated illness caused by Group A Streptococcus producing erythrogenic toxin. Presents with fever, pharyngitis, and a characteristic diffuse, erythematous, sandpaper-textured rash that blanches with pressure. Circumoral pallor, Pastia lines (linear petechiae in skin folds), and strawberry tongue are classic findings. Treatment: penicillin or amoxicillin for 10 days. Prevents rheumatic fever complications.

\medterm{Shaken Baby Syndrome}-- A form of abusive head trauma caused by violent shaking, typically in infants under 1 year. The triad: subdural hemorrhage, retinal hemorrhages, and encephalopathy. May present with irritability, lethargy, vomiting, seizures, and apnea. CT shows subdural hematoma; MRI shows diffuse axonal injury. Ophthalmologic exam reveals retinal hemorrhages. High mortality and long-term neurodevelopmental impairment.

\medterm{Sickle Cell Anemia}-- The most severe form of sickle cell disease (HbSS), causing chronic hemolytic anemia and vaso-occlusive crises. Features: jaundice, splenomegaly (infancy), pain crises, acute chest syndrome, stroke, and increased infection risk from functional asplenia. Diagnosis: hemoglobin electrophoresis showing HbS. Management: penicillin prophylaxis, hydroxyurea, pneumococcal vaccines, and folic acid supplementation.

\medterm{SIDS}-- See \hyperlink{Sudden Infant Death Syndrome}{Sudden Infant Death Syndrome}.

\medterm{Skeletal Survey}-- A series of radiographs obtained when child abuse is suspected. Includes AP and lateral views of the skull, chest (including oblique views), abdomen, spine, and bilateral upper and lower extremities. Identifies fractures of varying ages, which is highly suggestive of non-accidental trauma. Initial survey may be negative; repeat at 2 weeks may reveal healing fractures. Interpretation by pediatric radiologist recommended.

\medterm{Slipped Capital Femoral Epiphysis}-- Displacement of the femoral head relative to the femoral neck through the growth plate (physis) in adolescents. Most common hip disorder in adolescents (10-16 years). Risk factors: obesity, male sex, African American, and endocrine disorders (hypothyroidism, growth hormone deficiency). Presents with hip, groin, or knee pain with limping (antalgic gait). On examination: limited internal rotation and obligatory external rotation with hip flexion. Diagnosis: bilateral hip X-ray (AP and frog-leg lateral) shows femoral head displaced posteriorly and medially relative to the neck (Klein line). Treatment: surgical stabilization with in situ screw fixation of the affected and contralateral hip. Urgent: non-weight-bearing to prevent further slippage.

\medterm{Smith-Lemli-Opitz Syndrome}-- An autosomal recessive disorder of cholesterol metabolism caused by mutations in DHCR7, encoding 7-dehydrocholesterol reductase, the final enzyme in the cholesterol biosynthetic pathway. Features: microcephaly, characteristic facies (ptosis, anteverted nares, micrognathia, low-set ears), intellectual disability (mild to profound), syndactyly of the second and third toes (pathognomonic, 95\%), postaxial polydactyly, cleft palate, hypospadias (ambiguous genitalia in males), and congenital heart disease. Diagnosis: elevated serum 7-dehydrocholesterol levels and DHCR7 genetic testing. Treatment: dietary cholesterol supplementation (improves growth and behavior), and symptomatic management of associated anomalies. Prognosis correlates with severity of malformations and cholesterol levels.

\medterm{Speech Delay}-- A delay in language development compared to peers: by 12 months—no single words; 18 months—<10-20 words/no word combinations by 24 months; or regression. Causes: hearing impairment (must exclude—audiology), global developmental delay, autism spectrum disorder (with social communication deficits), expressive language disorder, psychosocial/environmental deprivation, and prematurity. Bilingualism does not cause delay. Evaluation: hearing test (essential), developmental assessment, speech-language evaluation, and screening for autism. Management: early intervention—speech therapy, parent training, treat underlying cause; monitor progress; red flags (loss of skills, poor response to name) prompt urgent referral.

\medterm{Spina Bifida}-- A neural tube defect: incomplete closure of the neural tube (days 22-28), causing spinal cord/meningeal malformation. Types: spina bifida occulta (hidden—no protrusion, often asymptomatic, hairy patch/hemangioma over defect), meningocele (meninges protrude), and myelomeningocele (cord + meninges protrude—most severe, neurologic deficits below the lesion: paralysis, sensory loss, bladder/bowel dysfunction, hydrocephalus, Chiari II malformation). Associated with folate deficiency; prevention: periconceptional folic acid. Prenatal: AFP (elevated), ultrasound. Management: multidisciplinary—surgical closure, VP shunt for hydrocephalus, urologic/orthopedic care, and rehabilitation.

\medterm{Spina Bifida Occulta}-- The mildest form of spina bifida, characterized by incomplete closure of the vertebral arch without herniation of neural elements. Often asymptomatic; may be detected by skin findings (dimple, hair tuft, lipoma) over the lower spine. Incidence approximately 1-2\% of the population. Usually requires no treatment unless associated with tethered cord or spinal cord symptoms.

\medterm{Stepping Reflex}-- A primitive reflex present from birth to 2 months, elicited by holding the infant upright with feet touching a flat surface. The infant makes stepping-like movements as if walking. Disappears by 2 months as voluntary motor control develops. Reappears around 9-10 months as voluntary walking emerges. Assessed during newborn and early infant neurologic examination.

\medterm{Stickler Syndrome}-- An autosomal dominant connective tissue disorder caused by mutations in collagen genes (COL2A1 most common, also COL11A1, COL11A2, COL9A1). Features: Pierre Robin sequence (micrognathia, glossoptosis, cleft palate, 30\%), high myopia (present in 90\% by 10 years), retinal detachment (highest risk of any pediatric syndrome), sensorineural hearing loss, and early-onset osteoarthritis with mild epiphyseal dysplasia. Pierre Robin sequence in Stickler syndrome mandates ophthalmologic screening. Diagnosis: clinical criteria (major: Pierre Robin-associated cleft palate, high myopia, retinal detachment, hearing loss, family history) and genetic testing. Management: ophthalmology — annual dilated fundus exams, retinal detachment prevention; audiology — hearing aids; orthopedics — joint protection and pain management; dental palatal repair; avoidance of contact sports.

\medterm{Strep Throat}-- Acute streptococcal pharyngitis caused by group A Streptococcus (S. pyogenes), most common in children 5-15 years. Features: sudden sore throat, fever, tonsillar exudate, palatal petechiae, anterior cervical lymphadenopathy, and headache/abdominal pain; cough, coryza, and oral ulcers suggest viral. Diagnosis: rapid antigen test or throat culture (confirmatory); clinical criteria (Centor/McIsaac). Complications: suppurative (peritonsillar abscess, otitis) and nonsuppurative (acute rheumatic fever, post-streptococcal glomerulonephritis). Management: antibiotics within 9 days to prevent rheumatic fever—penicillin (or amoxicillin), macrolide if allergic; analgesia and supportive care.

\medterm{Sturge-Weber Syndrome}-- A neurocutaneous syndrome (phakomatosis) characterized by a facial capillary malformation (port-wine stain) in the V1 distribution of the trigeminal nerve (forehead and upper eyelid) and ipsilateral leptomeningeal angiomatosis (venous malformations of the pia mater). Caused by a somatic activating mutation in GNAQ. Presents in infancy with the port-wine stain, which is the visible marker. Neurologic features: seizures (75-90\%, often focal or infantile spasms, typically starting in the first year), hemiparesis (contralateral to the angioma, 30-50\%), developmental delay, and intellectual disability. Ocular features: glaucoma (30-70\%, can be congenital or develop later, ipsilateral to the angioma), and choroidal hemangioma. Diagnosis: clinical with contrast-enhanced brain MRI (leptomeningeal enhancement, cerebral atrophy, gyriform calcification, and choroid plexus enlargement). Skull X-ray may show tram-track calcifications. Treatment: anticonvulsants for seizures; early hemispherotomy or corpus callosotomy for refractory epilepsy. Regular ophthalmologic monitoring for glaucoma (lifelong). Laser therapy for the port-wine stain. Surveillance: developmental assessments and neuroimaging.

\medterm{Subgaleal Hemorrhage}-- A life-threatening accumulation of blood in the subgaleal space (between the periosteum and the epicranial aponeurosis). Can contain up to 50\% of the infant's blood volume. Risk factors include vacuum-assisted delivery. Presents with rapidly enlarging, boggy scalp swelling that crosses suture lines, associated with anemia and hemodynamic instability. Requires urgent evaluation, blood transfusion, and resuscitation.

\medterm{Sucking Reflex}-- A primitive reflex present from birth, elicited by placing an object in the infant's mouth. The infant rhythmically sucks on the object. Coordinated sucking and swallowing typically develop by 34-36 weeks gestational age. Premature infants may lack coordinated suck-swallow, requiring gavage feeding. Absence in a term infant may indicate neurologic impairment. Important for successful breastfeeding and oral feeding.

\medterm{Sudden Infant Death Syndrome}-- SIDS: the sudden, unexplained death of an infant <1 year (peak 2-4 months) that remains unexplained after a thorough investigation (autopsy, scene examination, and review of clinical history). Risk factors: prone sleep, bed-sharing, soft bedding, overheating, maternal smoking, prematurity/low birth weight, and co-sleeping. Protective factors: supine sleep ('Back to Sleep'), firm mattress, room-sharing without bed-sharing, breastfeeding, pacifier, and no smoke exposure. Diagnosis is by exclusion after comprehensive investigation. Prevention: safe sleep education is the cornerstone. Grief support and counseling for families; differentiate from accidental suffocation and child abuse.

\medterm{Sudden Infant Death Syndrome (SIDS)}-- See \hyperlink{Sudden Infant Death Syndrome}{Sudden Infant Death Syndrome}.

\medterm{Surfactant}-- A phospholipid-protein complex produced by type II pneumocytes that reduces alveolar surface tension, preventing atelectasis. Deficiency causes respiratory distress syndrome in premature infants. Exogenous surfactant (beractant, poractant alfa) administered via endotracheal tube. Natural surfactant contains phosphatidylcholine, phosphatidylglycerol, and surfactant proteins A, B, C, and D. Production begins at approximately 24-28 weeks gestation.

\medterm{Syphilis}-- A bacterial STI caused by Treponema pallidum. Primary: chancre at inoculation site. Secondary: rash, mucous patches, condylomata lata. Latent: asymptomatic. Tertiary: gummas, cardiovascular, neurosyphilis. Diagnosis: non-treponemal tests (RPR/VDRL) confirmed by treponemal tests. Treatment: penicillin G benzathine IM. Congenital syphilis preventable with maternal screening and treatment.

\medterm{Tanner Staging}-- A scale for assessing sexual maturity during puberty. Stage 1: prepubertal. Stage 2: breast buds (girls), testicular enlargement (boys). Stage 3: breast enlargement, pubic hair. Stage 4: areolar development, more pubic hair. Stage 5: adult breasts and genitalia, adult hair pattern. Used to track pubertal progression and identify precocious or delayed puberty.

\medterm{Tdap Vaccine}-- Tetanus, Diphtheria, and acellular Pertussis booster given at 11-12 years and during each pregnancy (27-36 weeks gestation). Replaces one dose of Td. Important for pertussis protection, particularly for protecting newborns through maternal antibody transfer. Also given to adolescents and adults to maintain tetanus and diphtheria immunity.

\medterm{Teenage Pregnancy}-- Pregnancy in adolescents (usually 10-19 years), associated with higher maternal and fetal risks: preterm birth, low birth weight, preeclampsia, anemia, and increased neonatal mortality; also socioeconomic and educational consequences. Causes: unprotected intercourse, lack of contraception/education, poverty, and peer/family factors. Management: comprehensive antenatal care (nutrition, folate), psychosocial support, education continuation, and safe delivery; counseling on contraception postpartum to prevent recurrence. Prevention: sex education, access to contraception, and youth-friendly services. Multidisciplinary support (medical, social, educational) improves outcomes.

\medterm{Teething}-- The eruption of primary (deciduous) teeth, usually starting around 6 months (range 3-12), with lower central incisors first; ~20 teeth by 2-3 years. Features: irritability, drooling, gum swelling/tenderness, biting/chewing, disrupted sleep, mild low-grade fever; NOT a cause of high fever, diarrhea, or serious illness (which warrant evaluation). Management: teething rings, gum massage, cool (not frozen) items, acetaminophen/ibuprofen for discomfort; avoid topical benzocaine (risk). Emergence of permanent teeth continues through childhood (6-12 years). Differentiate from infections when fever is significant.

\medterm{Tetralogy of Fallot}-- The most common cyanotic congenital heart defect, accounting for 7-10\% of congenital cardiac anomalies. Characterized by four anatomical features: 1) Subpulmonary infundibular stenosis (right ventricular outflow tract obstruction — primary defect), 2) Ventricular septal defect (VSD), 3) Overriding aorta, and 4) Right ventricular hypertrophy (RVH). Associated with chromosome 22q11.2 deletion (DiGeorge syndrome). Clinical presentation: cyanosis ("blue baby") appearing in early infancy; hypercyanotic "tet spells" (triggered by crying, feeding, or exertion; managed by squatting or knee-to-chest position, which increases SVR and reduces right-to-left shunt). Physical exam: harsh holosystolic ejection murmur at upper left sternal border; single S2. Diagnostic imaging: chest X-ray shows classic "boot-shaped heart" (coeur en sabot) with decreased pulmonary vascular markings; echocardiogram confirms 4 defects. Treatment: surgical repair (VSD closure and RVOT relief) typically performed at 3-6 months of age.

\medterm{Thalassemia Major}-- Beta-thalassemia, an autosomal recessive hemoglobin disorder with absent or minimal beta-globin chain production. Presents in infancy with severe anemia, failure to thrive, hepatosplenomegaly, and skeletal deformities. Diagnosis by hemoglobin electrophoresis showing absent HbA and elevated HbF and HbA2. Treatment: chronic blood transfusions, iron chelation, and hematopoietic stem cell transplant.

\medterm{Tonic Neck Reflex}-- A primitive reflex present from birth to 6 months, also called fencer's pose. When the head is turned to one side, the arm on that side extends while the opposite arm flexes. Present in both directions. Asymmetry may indicate neurologic injury. Disappears as voluntary head control develops. One of several posture-related reflexes assessed in the newborn period.

\medterm{Toxoplasmosis}-- See infectious disease.

\medterm{Tracheoesophageal Fistula}-- An abnormal connection between the trachea and esophagus, often associated with esophageal atresia. Presents with recurrent aspiration, coughing with feedings, and cyanosis during meals. The H-type fistula may present later in infancy with chronic aspiration and failure to thrive. Diagnosis by esophagram or bronchoscopy. Treatment: surgical ligation and division of the fistula.

\medterm{Transfusion of Neonate}-- Administration of packed red blood cells to neonates for symptomatic anemia or hematocrit below transfusion threshold. Indications include acute blood loss, severe anemia (hematocrit less than 20\%), and respiratory distress. Volume: 10-20 mL/kg over 2-4 hours. Blood products must be CMV-negative, irradiated, and crossmatched. Complications include volume overload, infection, and transfusion reactions.

\medterm{Transposition of the Great Arteries}-- A cyanotic congenital heart defect where the aorta arises from the right ventricle and the pulmonary artery arises from the left ventricle (ventriculoarterial discordance). The result: two parallel circulations (systemic and pulmonary) with no mixing. Survival depends on associated shunts (ASD, VSD, PDA) to allow mixing of oxygenated and deoxygenated blood. Presents in the first days of life with severe cyanosis, tachypnea, and profound hypoxemia. Classic chest X-ray: egg-on-side appearance (narrow superior mediastinum). Treatment: immediate PGE1 infusion to maintain PDA patency, balloon atrial septostomy (Rashkind procedure) to improve mixing, and ultimately arterial switch operation (Jatene procedure) within the first 2 weeks of life.

\medterm{Treacher Collins Syndrome}-- An autosomal dominant craniofacial disorder caused by mutations in TCOF1 (most common), POLR1C, or POLR1D genes, encoding ribosomal RNA transcription factors. Bilateral, symmetric craniofacial anomalies: zygomatic hypoplasia (malar flattening), micrognathia, downward-slanting palpebral fissures, coloboma of the lower eyelid, microtia (small, malformed pinna), and conductive hearing loss (from middle ear anomalies). Intelligence is typically normal. Diagnosis: clinical and genetic testing. Management: multidisciplinary — airway assessment (possible tracheostomy for severe micrognathia), hearing aids or bone-anchored hearing aids for hearing loss, and staged reconstructive surgery (zygomatic and mandibular distraction, eyelid reconstruction, and external ear reconstruction).

\medterm{Tricuspid Atresia}-- Absence of the tricuspid valve, resulting in no direct communication between the right atrium and right ventricle. Blood flow depends on an atrial septal defect and ventricular septal defect for mixing. Presents with cyanosis. Staged surgical repair: modified Blalock-Taussig shunt in infancy, followed by Fontan procedure in early childhood. Oxygen and prostaglandin E1 used for stabilization.

\medterm{Trisomy 21}-- Down syndrome: the most common chromosomal disorder, caused by an extra chromosome 21 (trisomy 21; ~95 percent meiotic nondisjunction, mosaicism, translocation). Features: intellectual disability, characteristic facial features (flat facies, upslanting palpebral fissures, epicanthic folds, macroglossia), hypotonia, congenital heart disease (~50 percent—AVSD, VSD), GI anomalies (duodenal atresia), hypothyroidism, hearing/vision problems, and increased risk of leukemia and Alzheimer disease. Antenatal screening (NIPT, quad test) and diagnosis (amniocentesis/CVS). Management: multidisciplinary—cardiac repair, thyroid monitoring, audiology/ophthalmology, early intervention, education, and adult care; life expectancy now ~60 years.

\medterm{Truncus Arteriosus}-- A congenital heart defect where a single large vessel arises from both ventricles, receiving a mixed blood supply. Associated with a large ventricular septal defect and abnormal semilunar valve. Presents with cyanosis and heart failure early in life. Requires surgical repair in the neonatal period, involving VSD closure and placement of a right ventricle-to-pulmonary artery conduit.

\medterm{Tuberculin Skin Test}-- Intradermal injection of purified protein derivative (PPD) to test for Mycobacterium tuberculosis infection. Read at 48-72 hours measuring induration (not erythema). Positive: greater than 5 mm in HIV or close contacts; greater than 10 mm in high-risk groups; greater than 15 mm in low-risk populations. Does not differentiate latent from active infection. False positives from BCG vaccination.

\medterm{Tuberculosis in Children}-- Infection with Mycobacterium tuberculosis, often from household contacts. Latent TB infection (LTBI) positive tuberculin skin test or interferon-gamma release assay with no symptoms or radiographic findings. Active TB presents with cough, fever, weight loss, and lymphadenopathy. Diagnosis by PPD, chest X-ray, and sputum culture. Treatment: isoniazid, rifampin, pyrazinamide, and ethambutol for active disease.

\medterm{Tuberous Sclerosis Complex}-- An autosomal dominant genetic disorder caused by mutations in TSC1 (hamartin) or TSC2 (tuberin) genes, leading to mTOR pathway activation. See neurology.

\medterm{Tympanostomy Tubes}-- Small tubes inserted through the tympanic membrane to ventilate the middle ear and prevent fluid accumulation. Indicated for recurrent acute otitis media (3 episodes in 6 months or 4 in 12 months) or persistent otitis media with effusion greater than 3 months with hearing loss. Performed as outpatient procedure under general anesthesia. Most tubes extrude spontaneously within 6-18 months.

\medterm{Tyrosinemia}-- An autosomal recessive inborn error of tyrosine metabolism. Type I (fumarylacetoacetate hydrolase deficiency, FAH mutations) is the most severe: presents in infancy with liver failure, renal tubular dysfunction (Fanconi syndrome — hypophosphatemic rickets), neurologic crises (painful paresthesias, autonomic instability, porphyria-like symptoms from succinylacetone accumulation), and hepatocellular carcinoma risk. Type II (tyrosine aminotransferase deficiency, TAT mutations): oculocutaneous tyrosinemia with painful palmar/plantar hyperkeratosis and corneal ulcers; no liver or kidney disease. Type III (4-hydroxyphenylpyruvate dioxygenase deficiency, HPD mutations): rare, mild intellectual disability and ataxia. Newborn screening shows elevated tyrosine and succinylacetone (pathognomonic for Type I). Diagnosis: urine succinylacetone (Type I), plasma amino acids, and FAH/TAT/HPD genetic testing. Treatment: Type I — nitisinone (NTBC, 2-(2-nitro-4-trifluoromethylbenzoyl)-1,3-cyclohexanedione) blocks tyrosine catabolism upstream of FAH, preventing toxic metabolite accumulation, combined with a tyrosine- and phenylalanine-restricted diet. Liver transplantation for refractory liver disease or HCC. Type II — dietary restriction of tyrosine and phenylalanine.

\medterm{Undescended Testis}-- Cryptorchidism: failure of one or both testes to descend into the scrotum, the most common congenital male genital anomaly (preterm 30 percent, term 3-4 percent); often descends spontaneously in the first months. Complications: impaired fertility, increased risk of testicular cancer, torsion, and hernia. Diagnosis: clinical (absent testis in the scrotum—exclude retractile); imaging rarely needed; bilaterally impalpable—hormone/ultrasound evaluation. Management: surgical orchidopexy recommended at 6-12 (ideally 6-18) months to preserve fertility and allow surveillance; if bilateral and hypogonadism—hormonal evaluation. Monitor for malignancy with regular examination after surgery.

\medterm{Urea Cycle Disorder}-- A group of inborn errors of ammonia metabolism caused by deficiency of any of the six enzymes or two transporters in the urea cycle. X-linked: ornithine transcarbamylase deficiency (OTC, most common, 50\%). Autosomal recessive: N-acetylglutamate synthase (NAGS), carbamoyl phosphate synthetase I (CPS1), argininosuccinate synthetase (citrullinemia), argininosuccinate lyase (argininosuccinic aciduria), and arginase 1 (argininemia). Presents with hyperammonemic encephalopathy: poor feeding, vomiting, lethargy, tachypnea/respiratory alkalosis (from central hyperventilation), ataxia, seizures, and coma. Onset varies from neonatal (complete deficiency) to late-onset (partial deficiency). Diagnosis: plasma ammonia (elevated >150 \(\mu\)mol/L in neonates is symptomatic), quantitative plasma amino acids (citrulline, arginine, glutamine, alanine), urine orotic acid, and genetic testing. Treatment: acute hyperammonemia — IV sodium phenylacetate/sodium benzoate (Ammonul), IV arginine, nitrogen-free IV fluids, and hemodialysis for severe (>500 \(\mu\)mol/L). Chronic management — low-protein diet, ammonia-scavenging drugs (sodium phenylbutyrate, glycerol phenylbutyrate, arginine or citrulline supplementation), and liver transplantation for severe OTC deficiency. NAGS deficiency is uniquely treatable with oral N-carbamylglutamate.

\medterm{Vaccination Schedule}-- The recommended immunization schedule for infants, children, and adolescents (e.g., WHO/CDC): birth—HepB; 2 months—DTaP, IPV, Hib, PCV13, RV; 4 months—same; 6 months—HepB, IPV, PCV, RV, influenza (seasonal); 12 months—MMR, varicella, Hib, PCV; 18 months—DTaP, HepA; 4-6 years—DTaP, IPV, MMR, varicella; 9-16 years—HPV; 11-12 years—Tdap, meningococcal (MenACWY); annual influenza. Catch-up schedules for missed doses. Timing ensures protection at vulnerable ages; vaccines are safe and effective; document and counsel on side effects. School entry requirements often enforce immunization. See also \hyperlink{India Immunization Schedule}{India Immunization Schedule}.

\medterm{India Immunization Schedule}-- The Universal Immunization Schedule (UIP) of the Government of India, providing free vaccines to all children, adolescents, and pregnant women through the public health system. The childhood schedule: at birth---BCG (single dose, intradermal, upper arm, against tuberculosis), OPV-0 (zero dose, oral, against poliovirus), and Hepatitis B birth dose (injection); 6 weeks---OPV-1, Pentavalent-1 (DPT+HepB+Hib, injection), Rotavirus Vaccine-1 (oral), Pneumococcal Conjugate Vaccine-1 (PCV, injection), and fractional IPV-1 (fIPV, intradermal); 10 weeks---OPV-2, Pentavalent-2, Rotavirus-2; 14 weeks---OPV-3, Pentavalent-3, Rotavirus-3, PCV-2, and fIPV-2; 9-12 months---Measles-Rubella-1 (MR-1, injection), Japanese Encephalitis-1 (JE-1, in endemic districts, injection), PCV booster, and fIPV-3 (full intramuscular dose); 16-24 months---MR-2, JE-2, DPT booster-1, and OPV booster. The adolescent schedule: 5-6 years---DPT booster-2; 10 years---Td (Tetanus and adult Diphtheria); 16 years---Td booster. The maternal schedule: Td-1 as early as possible in pregnancy, Td-2 at least 4 weeks later (one booster dose sufficient if two Td doses received within the preceding 3 years). India uses the bivalent oral polio vaccine (bOPV, types 1 and 3) and fractional IPV (0.1 mL intradermal at 6 and 14 weeks) to optimise polio immunity while managing vaccine supply. The pentavalent vaccine reduces the number of injections from five to one at each visit. Japanese encephalitis vaccine is given only in districts with high JE burden. The schedule aligns with WHO recommendations and has contributed to India achieving polio-free certification in 2014.

\medterm{VACTERL Association}-- A non-random association of congenital anomalies, diagnosed when at least three of the following are present: Vertebral anomalies (hemivertebrae, fused vertebrae, scoliosis, 70\%), Anorectal malformations (imperforate anus, 80\%), Cardiac defects (VSD, tetralogy of Fallot, 50-80\%), Tracheo-Esophageal fistula (with or without esophageal atresia, 70\%), Renal anomalies (renal agenesis, horseshoe kidney, hydronephrosis, 60\%), and Limb defects (radial dysplasia, thumb anomalies, polydactyly, 50-70\%). The acronym also includes non-VACTERL findings: single umbilical artery and rib anomalies. The cause is unknown; most cases are sporadic. Diagnosis: clinical with systematic evaluation of involved organ systems. Management: staged surgical repair (imperforate anus colostomy, TEF repair, cardiac surgery), and long-term follow-up for renal function, growth, and development.

\medterm{Varicella Vaccine}-- Live attenuated vaccine given at 12-15 months and 4-6 years. Protects against chickenpox (varicella-zoster virus). Two doses provide greater than 90\% protection. Contraindicated in immunocompromised and pregnant individuals. Also reduces risk of herpes zoster later in life. Universal vaccination significantly reduced varicella incidence and complications.

\medterm{Viral Meningitis}-- Inflammation of the meninges caused by viruses, usually Enteroviruses. Presents similarly to bacterial meningitis but generally less severe. CSF shows lymphocytic pleocytosis, normal glucose, and mildly elevated protein. Self-limited in most cases; supportive treatment. Differentiated from bacterial meningitis by CSF analysis to prevent unnecessary prolonged antibiotic therapy.

\medterm{Vision Screening}-- Screening for visual problems in children: newborn—red reflex examination (cataract, retinoblastoma); 3-5 years—visual acuity, strabismus screening (cover test, photoscreening), and amblyopia risk factors. Goals: detect amblyopia ('lazy eye'), strabismus, and refractive errors early (before 5-7 years, the critical period), since treatment (patching, glasses) is most effective early. Risk factors: prematurity, family history, syndromes. Referral criteria: acuity below age norms, significant refractive error, strabismus, or abnormal red reflex. Regular screening at well-child visits and school entry; early referral to pediatric ophthalmology prevents permanent vision loss.

\medterm{Whooping Cough}-- A highly contagious bacterial infection caused by Bordetella pertussis, producing characteristic paroxysmal cough with whooping inspiratory sound. See \hyperlink{Pertussis}{Pertussis}.

\medterm{Williams Syndrome}-- A microdeletion syndrome involving the elastin gene on chromosome 7q11.23. Features: elfin facies, supravalvular aortic stenosis, hypercalcemia, intellectual disability with strong verbal skills, and overly friendly personality. Often associated with connective tissue abnormalities. Management: cardiac monitoring, calcium management, and educational support.

\medterm{Wilms' tumor (Nephroblastoma)}-- The most common renal malignancy in children, accounting for ~95\% of pediatric renal tumors (peak incidence 2-5 years). It is an embryonal tumor arising from metanephric blastemal cells. Associated syndromes: WAGR (Wilms, Aniridia, Genitourinary anomalies, Intellectual Disability (formerly mental retardation) — WT1 deletion on 11p13), Beckwith-Wiedemann syndrome (hemihypertrophy, macroglossia, omphalocele — WT2 locus 11p15), and Denys-Drash syndrome (WT1 mutation, male pseudohermaphroditism, nephrotic syndrome). Pathology: triphasic histology (blastemal, stromal, epithelial components). Anaplastic histology (focal or diffuse) confers a worse prognosis. Presents with a large, asymptomatic abdominal mass (often noted by parent while bathing the child), abdominal pain, haematuria (microscopic or gross), and hypertension (from renin secretion by the tumor or renal artery compression). Workup: abdominal ultrasound (distinguishes solid tumor from hydronephrosis, cystic lesions, neuroblastoma) and CT abdomen (staging: tumor extent, lymph node involvement, bilateral disease, IVC thrombus). Staging: Stage I (limited to kidney, capsule intact, completely resected), Stage II (extends beyond kidney, completely resected), Stage III (gross residual disease, positive lymph nodes, tumor spillage), Stage IV (haematogenous metastases — lung, liver, bone, brain), Stage V (bilateral renal involvement). Treatment: multimodal therapy based on staging and histology. Surgery: radical nephrectomy (transperitoneal approach, remove kidney + Gerota fascia + ipsilateral ureter + lymph node sampling). Chemotherapy: regimen depends on stage and histology (favourable vs. anaplastic) — vincristine, dactinomycin, doxorubicin. Radiation: for high-stage disease (Stage III-IV, anaplastic histology). Bilateral Wilms: neoadjuvant chemotherapy followed by nephron-sparing surgery. Prognosis: overall cure rate exceeds 90\% (favourable histology). Anaplastic histology and metastatic disease have lower survival rates (50-70\%). Long-term follow-up for renal function, cardiac function (doxorubicin), and second malignancies.

\medterm{Corrected Age}-- The age of a preterm infant calculated by subtracting the number of weeks born before 40 weeks of gestation from chronological age. Corrected age helps interpret growth and developmental milestones during early childhood, commonly until about 2 years of age.

\medterm{Developmental Surveillance}-- The ongoing process of eliciting caregiver concerns, observing the child, documenting developmental history, and monitoring milestones during routine care. It differs from formal developmental screening, which uses standardized instruments at specified ages.

\medterm{Developmental Regression}-- Loss of previously acquired developmental skills. Regression is a neurologic red flag requiring prompt evaluation for genetic, metabolic, neurodegenerative, epileptic, infectious, or structural causes.

\medterm{Precocious Puberty}-- Onset of secondary sexual characteristics earlier than the accepted age threshold, commonly before 8 years in girls or 9 years in boys. It may be central from early activation of the hypothalamic-pituitary-gonadal axis or peripheral from excess sex steroids.

\medterm{Delayed Puberty}-- Failure to begin or progress through puberty by the expected age, such as absence of breast development by 13 years in girls or testicular enlargement by 14 years in boys. Causes include constitutional delay, chronic disease, gonadal failure, and hypothalamic or pituitary disorders.

\medterm{Dehydration in Children}-- A deficit of total body water, usually caused by gastrointestinal losses, inadequate intake, fever, or renal losses. Clinical assessment includes mental status, heart rate, mucous membranes, tears, capillary refill, urine output, weight change, and perfusion rather than a single sign.

\medterm{Pediatric Shock}-- Inadequate tissue oxygen delivery or utilization in a child, most often caused by hypovolemia, distributive disease, cardiogenic dysfunction, or obstruction. Children may maintain blood pressure until late; tachycardia, abnormal perfusion, altered mental status, and oliguria are important early warning signs.

\medterm{Intussusception}-- Telescoping of one segment of intestine into an adjacent segment, causing intermittent obstruction and impaired venous drainage. It commonly presents in infants with episodic colicky pain, vomiting, lethargy, and sometimes currant-jelly stool; ultrasound and pneumatic or hydrostatic reduction are standard approaches when appropriate.

\medterm{Hypertrophic Pyloric Stenosis}-- Hypertrophy and hyperplasia of the pyloric muscle causing gastric outlet obstruction, typically in young infants. Progressive nonbilious projectile vomiting, hunger after vomiting, dehydration, and hypochloremic hypokalemic metabolic alkalosis may occur.

\medterm{Foreign-Body Aspiration}-- Entry of an object into the larynx, trachea, or bronchi, most often in toddlers. Sudden cough, choking, unilateral wheeze, or focal air trapping may occur; complete airway obstruction is an immediate emergency and requires age-appropriate resuscitation.

\medterm{Pediatric Appendicitis}-- Acute inflammation of the vermiform appendix, usually caused by luminal obstruction. Children may have evolving periumbilical-to-right-lower-quadrant pain, anorexia, vomiting, fever, and guarding, but presentation can be atypical in young children; timely surgical and antimicrobial assessment reduces perforation risk.

\medterm{Congenital Hypothyroidism}-- Insufficient thyroid hormone present at birth, caused by thyroid dysgenesis, dyshormonogenesis, or central disease. Early screening and treatment with levothyroxine prevent severe neurodevelopmental impairment.

\medterm{Phenylketonuria (PKU)}-- An inherited disorder, usually due to phenylalanine hydroxylase deficiency, causing accumulation of phenylalanine and reduced tyrosine availability. Untreated disease can cause intellectual disability, seizures, eczema, and a musty odor; newborn screening and dietary treatment are effective.

\medterm{Neonatal Hypoxic-Ischemic Encephalopathy (HIE)}-- Brain dysfunction in a newborn caused by inadequate oxygenation and/or cerebral blood flow around the time of birth. Moderate or severe HIE may produce altered consciousness, abnormal tone, seizures, and multiorgan injury; selected term or near-term infants may benefit from early therapeutic hypothermia.

\medterm{Phototherapy}-- Treatment of neonatal unconjugated hyperbilirubinemia using light wavelengths that convert bilirubin into water-soluble photoisomers excreted without hepatic conjugation. Indications depend on gestational age, postnatal age, bilirubin level, and neurotoxicity risk factors.

\medterm{Catch-Up Immunization}-- Vaccination of a person who is behind the recommended schedule using age- and risk-appropriate intervals. Doses generally do not need to be restarted because of delay, but minimum ages and intervals must be respected.

\medterm{Sentinel Injury}-- A minor but potentially significant injury in a nonmobile infant, such as unexplained bruising, oral injury, or a fracture, that may precede more severe abuse. It warrants careful history, complete examination, safeguarding assessment, and documentation.

\medterm{Abusive Head Trauma}-- Injury to an infant or child caused by inflicted blunt force, shaking, impact, or a combination of mechanisms. Findings may include subdural hemorrhage, retinal hemorrhages, encephalopathy, seizures, and fractures; suspected cases require urgent medical stabilization and safeguarding evaluation.

\medterm{Sudden Unexpected Infant Death (SUID)}-- The sudden and unexpected death of an infant younger than 1 year, whether or not the cause is known before investigation. Sudden infant death syndrome is the unexplained subtype after complete scene, history, and autopsy evaluation.

\medterm{Fencer's Pose}-- See \hyperlink{Tonic Neck Reflex}{Tonic Neck Reflex}.

\medterm{Pediculosis Capitis}-- See \hyperlink{Head lice}{Head lice}.

\medterm{Growth Retardation}-- Abnormally slow physical growth or failure to reach expected height or weight for age and developmental stage. Causes include inadequate nutrition, chronic systemic disease, endocrine or genetic disorders, malabsorption, psychosocial deprivation, and constitutional variation; assessment uses growth charts and clinical context.

\medterm{Enuresis}-- Repeated involuntary urination, usually during sleep in a child who is developmentally expected to remain dry. It may be primary or secondary and requires assessment for constipation, urinary infection, diabetes, sleep or developmental disorders, psychosocial stressors, and lower-urinary-tract dysfunction when clinically indicated.

\medterm{Iniencephaly}-- A rare, usually lethal neural-tube defect characterized by occipital bone defect, severe retroflexion of the head, and marked shortening or distortion of the cervical and thoracic spine. It is often identified prenatally and may occur with other congenital anomalies.

\medterm{Otocephaly}-- A severe congenital craniofacial malformation characterized by absent or markedly underdeveloped mandible, ventral displacement or fusion of the ears, and associated facial and airway abnormalities. It is usually lethal because of profound impairment of breathing and feeding.

\medterm{Bedwetting}-- Involuntary urination during sleep after the age at which nighttime bladder control is usually established. It may be primary or secondary to urinary infection, constipation, diabetes, sleep disorders, stress, or other medical conditions and should be assessed when persistent, new, or accompanied by daytime symptoms.
